% Editorial revision of the completed-study manuscript at 477304f.
% Scientific evidence and full existing proof are preserved; see README.md.
% No LaTeX compilation, PDF generation, rendering, or visual validation performed.
% Figures are byte-identical original SVGs. Future compilation requires the
% standard svg package and Inkscape; authorised shell escape may be required.
% Evidence hyperlinks resolve relative to this directory in a repository checkout.
\documentclass[11pt,a4paper]{article}
\usepackage[T1]{fontenc}
\usepackage[utf8]{inputenc}
\usepackage{lmodern}
\usepackage{textcomp}
\usepackage[margin=25mm]{geometry}
\usepackage{amsmath,amssymb,amsthm}
\usepackage{graphicx,svg}
\usepackage{array,booktabs,longtable,tabularx,xltabular}
\usepackage{caption}
\usepackage{algorithm}
\usepackage{fvextra}
\usepackage{xurl}
\usepackage[hidelinks]{hyperref}
\newtheorem{proposition}{Proposition}
\newcolumntype{Y}{>{\raggedright\arraybackslash}X}
\setcounter{secnumdepth}{2}
\newcommand{\appendixsection}[2]{%
  \refstepcounter{section}%
  \section*{Appendix \thesection. #2}%
  \addcontentsline{toc}{section}{Appendix \thesection. #2}%
  \label{#1}%
  \setcounter{table}{0}%
}
\title{TrafficTwin: Admission and Dispatch Semantics in Vehicular Edge Computing under a Frozen MAPPO Policy}
\author{}
\date{}
\begin{document}
\maketitle
% BEGIN SOURCE BLOCK 002 paragraph
S M Abdulla Al Mamun\\
The University of Manchester\\
Master's dissertation --- for author review, 9 September 2026\\
Supervisor: Dr Sandra Sampaio
% END SOURCE BLOCK 002

% BEGIN SOURCE BLOCK 003 heading
\section*{Abstract}\phantomsection\label{sec:abstract}
% END SOURCE BLOCK 003

% BEGIN SOURCE BLOCK 004 paragraph
A vehicle's strongest radio link can lead to a busy computing server, and admitting its task does not guarantee a result before the deadline. TrafficTwin investigates how infrastructure scheduling changes offered-task deadline attainment under a fixed, pretrained vehicle offloading policy. The comparison separates radio ingress, execution placement, admission and workload reservation.
% END SOURCE BLOCK 004

% BEGIN SOURCE BLOCK 005 paragraph
Common-target least-workload placement underperformed execution at ingress, while causal per-task placement improved it. This reversal appeared in an adaptive incident study and a separate four-draw morning replication that excluded its inspected pilot. A later confirmation completed 32 evaluations across eight fresh joint fleet/evaluator-seed blocks. Per-task minus ingress was +4.137 percentage points; common-target minus ingress was \ensuremath{-}3.504. Per-task also exceeded causal cyclic spreading by +0.631 points, with a Bonferroni simultaneous 95\% interval of [+0.511, +0.750]. All actions matched, although feedback was permitted and small observation/logit differences occurred in one block.
% END SOURCE BLOCK 005

% BEGIN SOURCE BLOCK 006 paragraph
Retrospective analysis explains recurring common destinations through shared assignment, queue clearing and deterministic ties. The exact-model proof establishes when fixed-target reconciliation reproduces sequential admission; rounding limits its application to executable arithmetic. Original morning task accounting locates most gains in reduced gate rejection, while retaining individual losses. Separate CPU fixtures show that workload-aware targeting costs more than cyclic spreading, despite outperforming the common-target implementation in scheduler-only timing.
% END SOURCE BLOCK 006

% BEGIN SOURCE BLOCK 007 paragraph
The findings support a bounded benefit from workload awareness beyond spreading. One actor, the same morning trace for confirmation, global service-work information and aggregate timing restrict generalisation and physical interpretation. Historical capacity findings separately retain inconclusive statistics and unresolved scoring uncertainty; the later gate-enabled results do not resolve that gap.
% END SOURCE BLOCK 007

% BEGIN SOURCE BLOCK 008 heading
\section{Introduction}\label{sec:1}
% END SOURCE BLOCK 008

% BEGIN SOURCE BLOCK 009 heading
\subsection{Problem and motivation}\label{sec:1.1}
% END SOURCE BLOCK 009

% BEGIN SOURCE BLOCK 010 paragraph
Vehicular edge computing concerns where a vehicle's computational tasks should run and whether their results can return before a deadline. A task may execute locally, use a neighbouring vehicle through vehicle-to-vehicle communication (V2V), or enter roadside infrastructure through vehicle-to-infrastructure communication (V2I). A roadside unit (RSU) combines a radio access point with a simulated computing server in this study. Offloading therefore couples communication conditions with the work already waiting for computation \cite{ref1}.
% END SOURCE BLOCK 010

% BEGIN SOURCE BLOCK 011 paragraph
The \textbf{ingress RSU} is the roadside unit through which a vehicle reaches the infrastructure. TrafficTwin chooses ingress using the strongest simulated radio link. \textbf{Execution placement} is the subsequent choice of server for the task. Keeping execution at ingress avoids logical forwarding, but strong radio quality says nothing about that server's backlog. Assigning the work to another RSU can reduce waiting while introducing a different information and coordination requirement. Figure~\ref{fig:1} shows the decision boundary studied here: the vehicle selects an offloading mode, while the infrastructure resolves the execution destination and admission \hyperref[source-s4]{[S4]}.
% END SOURCE BLOCK 011

% BEGIN SOURCE BLOCK 012 paragraph
Admission is itself distinct from timely completion. A finite waiting room can accept another task even when the queued service would make it late. TrafficTwin's deadline-aware gate improves on a count-only limit by inspecting backlog, but it does not include the candidate's own service or transmission time. An admitted task can consequently miss its deadline. Evaluation must retain rejected tasks in the offered population; measuring success only among admitted work would allow selective rejection to change the population being compared.
% END SOURCE BLOCK 012

% BEGIN SOURCE BLOCK 013 paragraph
The supplied multi-agent proximal policy optimisation (MAPPO) actor provides a fixed upstream policy for isolating infrastructure interventions. Its weights are held constant, removing retraining as an explanation for a changed ranking, while matched inputs and recorded action checks distinguish controlled variation from possible feedback through energy and observations. This design studies the scheduling implementation supporting an existing actor rather than seeking a new learned controller.
% END SOURCE BLOCK 013

% BEGIN SOURCE BLOCK 014 paragraph
The research began with accounting and capacity. An apparent load-balancing benefit depended on the admission comparator; two least-workload implementations then produced opposite rankings against ingress. Explaining target timing, visible reservations and admission outcomes makes this comparison assessable beyond scheduler names or workload balance \hyperref[source-s0]{[S0]}--\hyperref[source-s4]{[S4]}.
% END SOURCE BLOCK 014

% BEGIN SOURCE BLOCK 015 heading
\subsection{Related work and the position of this study}\label{sec:1.2}
% END SOURCE BLOCK 015

% BEGIN SOURCE BLOCK 016 paragraph
The closest intellectual comparison is task dispatch to parallel servers, particularly the information used to estimate waiting time and the point at which an assignment changes that information. RL supplies the upstream policy, but it does not define the downstream scheduling problem. The following synthesis positions the completed experiments retrospectively: these papers were newly inspected during manuscript revision and are not presented as having guided the historical experiments.
% END SOURCE BLOCK 016

% BEGIN SOURCE BLOCK 017 paragraph
\textbf{Queue count and remaining work.} Shortest-queue dispatch compares numbers of jobs; shortest-workload dispatch compares their outstanding service requirements. With heterogeneous task sizes, these orderings can disagree. Harchol-Balter, Crovella and Murta explicitly study dynamic least-work-remaining assignment when service demand is known, alongside random, round-robin and size-based assignment. Their model assumes immediate per-arrival assignment, identical hosts, non-preemptive first-come-first-served service and Poisson arrivals. Their analysis shows that preferred assignment policies depend on task-size variability \cite{ref2}. Consequently, neither choosing the least workload nor achieving balanced allocation establishes universal optimality. TrafficTwin's inherited \texttt{jsq} name is misleading if read literally: the implemented selector compares service milliseconds, while a separate task count enforces finite admission capacity.
% END SOURCE BLOCK 017

% BEGIN SOURCE BLOCK 018 paragraph
\textbf{Batch decisions and reservations.} Batching does not inherently imply sending a batch to one server. Sparrow pools probes across a parallel job and spreads its tasks over selected workers; late binding queues reservations and assigns tasks when workers become ready. Its discussion identifies both queue-count prediction errors and races between schedulers using apparently idle workers \cite{ref9}. This distinguishes two issues that a generic ``batch versus per-task'' description would conflate: sharing information across candidates can help, whereas committing many assignments against one unchanged destination can concentrate work. TrafficTwin evaluates the latter implementation. Its immediate service-work reservation is also different from Sparrow's worker-triggered task binding: TrafficTwin assumes the necessary information and acknowledgement are already available.
% END SOURCE BLOCK 018

% BEGIN SOURCE BLOCK 019 paragraph
Ying, Srikant and Kang make the reservation distinction especially explicit: batch-filling samples queue counts, assigns tasks one by one and updates the chosen queue after each assignment. Their analysis assumes identical servers, exponential service, Poisson batch arrivals and random ties \cite{ref13}. This is close prior art for sequential within-batch updates. Although the evaluator comments mention their batch-filling work, its common-target reconciliation does not implement that destination rule. TrafficTwin's causal change is therefore an adaptation of established dispatch semantics, evaluated with workload information and deadline admission, rather than a novel batching principle.
% END SOURCE BLOCK 019

% BEGIN SOURCE BLOCK 020 paragraph
\textbf{Local information and common destinations.} Vargaftik, Keslassy and Orda's Local Shortest Queue (LSQ) policies maintain possibly outdated queue-count estimates at multiple dispatchers. Their slotted model can send all jobs arriving at one dispatcher in a slot to one locally shortest queue, with random tie-breaking. Algorithms update local estimates by the jobs sent and refresh selected entries through communication. Their stability result requires stated arrival/service assumptions and bounded expected estimation error \cite{ref10}. Common destinations and local updates therefore predate TrafficTwin. Here they use service-work information, deterministic ties, a backlog gate, five ordered slots and one-second aggregate draining. Neither the LSQ stability theorem nor its distributed communication results transfer directly to these finite-horizon deadline outcomes.
% END SOURCE BLOCK 020

% BEGIN SOURCE BLOCK 021 paragraph
\textbf{Admission and useful completion.} Niño-Mora jointly studies admission and routing to heterogeneous multiserver queues, charging separately for rejection and admitted deadline misses. The model uses Poisson arrivals, exponential service, soft deadlines and state-dependent index policies; admitted late jobs remain until completion \cite{ref11}. This establishes admission as an optimisation decision with consequences beyond the admitted population. TrafficTwin instead preserves a simple inherited backlog filter to make the placement contrast interpretable. The gate does not estimate the complete response time or promise success. Counting successes over all offered tasks makes a reject-all policy score zero and prevents selective admission from concealing failures. That denominator is a deliberate measurement choice, not a newly invented admission algorithm.
% END SOURCE BLOCK 021

% BEGIN SOURCE BLOCK 022 paragraph
Table~\ref{tab:1} locates the comparison along operational dimensions. The papers identify relevant mechanisms and alternatives; they are not external benchmark arms. Their arrival processes, objectives and resource models differ too much for published performance percentages to be compared as controlled results.
% END SOURCE BLOCK 022

% BEGIN SOURCE BLOCK 023 table
\begingroup
\small
\begin{xltabular}{\textwidth}{@{}YYYY@{}}
\caption{Closest-work comparison. ``Work'' means remaining service demand; ``count'' means queued jobs. TrafficTwin rows describe the frozen implementation, not a scheduling-family definition.}\label{tab:1} \\
\toprule
Work / decision layer and arrivals & Information and update semantics & Admission / objective & Relationship to TrafficTwin \\
\midrule\endfirsthead
\toprule
Work / decision layer and arrivals & Information and update semantics & Admission / objective & Relationship to TrafficTwin \\
\midrule\endhead
\bottomrule\endfoot
Harchol-Balter et al. \cite{ref2}: dispatch each arrival & Known service work; central per-arrival choice & Mean waiting time; size-normalised waiting & Closest least-workload comparator; no vehicular gate \\
Sparrow \cite{ref9}: parallel-job task placement & Sampled workers; queued reservations; worker replies bind tasks & Job response time; placement constraints & Batch spreading differs from one common destination; communication is explicit \\
Ying et al. \cite{ref13}: tasks within an arriving batch & Sampled counts; update after each assignment; random ties & Delay and sampling cost & Closest within-batch update rule; no deadline gate \\
LSQ \cite{ref10}: dispatcher batch per time slot & Local queue counts; own-assignment increments and sampled/pushed updates & Stability under specified assumptions & Common routing is established; random ties and distributed estimates differ \\
Niño-Mora \cite{ref11}: joint admission/routing per arrival & Queue counts and model-based indices & Rejection and deadline-miss costs & Supports joint evaluation, not the inherited backlog predicate \\
TrafficTwin common-target: one choice per substep & Global service work; vectorised candidate offsets and three reconciliations & Backlog gate plus finite task limit & Frozen vehicle policy; instantaneous logical forwarding \\
TrafficTwin per-task: ordered candidate scan & Global service work; actual admitted work reserved immediately & Same predicate and limit, causal admission & Tested implementation change; global information is assumed \\
\end{xltabular}
\endgroup
% END SOURCE BLOCK 023

% BEGIN SOURCE BLOCK 024 paragraph
\textbf{Attribution and simulation validity.} Sargent distinguishes verification of a computer model from operational validation for its intended application \cite{ref12}. TrafficTwin's conservation and replay checks primarily establish the former; observed physical RSU service and network measurements would be needed for stronger deployment claims. SUMO supplies a microscopic simulation framework \cite{ref8}, but importing its trajectories does not independently validate a coupled radio/computing model. This distinction explains why a source-correct result can be scientifically useful while remaining conditional on simulation semantics.
% END SOURCE BLOCK 024

% BEGIN SOURCE BLOCK 025 paragraph
PPO and MAPPO identify the reused actor's provenance \cite{ref3}, \cite{ref4}, without establishing checkpoint optimality or identical actions under changed observations. Henderson et al., Agarwal et al. and Gorsane et al. motivate exposing implementation choices and finite-run uncertainty \cite{ref5}, \cite{ref6}, \cite{ref7}. TrafficTwin checks actions and uses draws or joint-seed blocks as replications; these references do not establish its small-sample distributional assumption.
% END SOURCE BLOCK 025

% BEGIN SOURCE BLOCK 026 heading
\subsection{Aim, research questions and contribution}\label{sec:1.3}
% END SOURCE BLOCK 026

% BEGIN SOURCE BLOCK 027 paragraph
The aim is to assess deterministic infrastructure scheduling under a frozen vehicle policy and explain its observed deadline outcomes. Questions developed through the programme. Morning replication and confirmation contrasts were declared before their new outcomes; mechanism analysis was post-hoc.
% END SOURCE BLOCK 027

% BEGIN SOURCE BLOCK 028 paragraph
\textbf{RQ1 --- Measurement and capacity.} What do the recorded accounting checks establish, and does increasing the RSU waiting-room limit at fixed compute service improve the archived score? Interpretation as admission-consistent deadline attainment additionally depends on the unresolved legacy scoring-mask exception (\S\,\ref{sec:3.1}). Sampling and measurement uncertainty are separate questions.
% END SOURCE BLOCK 028

% BEGIN SOURCE BLOCK 029 paragraph
\textbf{RQ2 --- Admission and dispatch semantics.} Under matched inputs and the inherited workload-based admission rule, how do ingress execution, common-target placement and causal per-task placement compare in the incident scenario? What does subsequent analysis establish about the operational mechanisms behind their different rankings? The explanatory analysis is retrospective, distinct from the original comparisons.
% END SOURCE BLOCK 029

% BEGIN SOURCE BLOCK 030 paragraph
\textbf{RQ3 --- Replication and workload awareness.} Does the reversal recur first across new morning fleet draws and then under fresh joint fleet/evaluator randomness? Under the same causal admission procedure, does workload-aware per-task targeting improve offered-task deadline attainment over cyclic spreading? The original four-draw replication and later eight-block confirmation answer separate stages of this question; their samples and interval families remain separate.
% END SOURCE BLOCK 030

% BEGIN SOURCE BLOCK 031 paragraph
\textbf{RQ4 --- Bounded sensitivities.} Within one incident fleet draw, how do aged placement-workload reports and a fixed forwarding overhead affect the comparison? These studies test specific model assumptions descriptively, without establishing population-level robustness or a physical backhaul design.
% END SOURCE BLOCK 031

% BEGIN SOURCE BLOCK 032 paragraph
The research contributes a measurement-qualified comparison of infrastructure implementations, separate replication and confirmation of the reversal, and an exact-model explanation bounded by numerical and information assumptions. Retrospective task-type and scheduler-cost analyses connect attainment to task outcomes and engineering trade-offs. E0/E1 provide preliminary measurement and scoping evidence; the final questions were developed through the programme, not fixed together at its outset. The research implementation and documented contributions are assessed in \S\,\ref{sec:3.7} \hyperref[source-s12]{[S12]}--\hyperref[source-s16]{[S16]}.
% END SOURCE BLOCK 032

% BEGIN SOURCE BLOCK 033 heading
\subsection{Scope and report structure}\label{sec:1.4}
% END SOURCE BLOCK 033

% BEGIN SOURCE BLOCK 034 paragraph
One actor, a provisional fleet model and two Manchester traces bound the study. Section~\ref{sec:2} specifies the model and design; Section~\ref{sec:3} evaluates findings and achievement; Section~\ref{sec:4} answers the questions. Appendices support verification.
% END SOURCE BLOCK 034

% BEGIN SOURCE BLOCK 035 heading
\section{Methodology}\label{sec:2}
% END SOURCE BLOCK 035

% BEGIN SOURCE BLOCK 036 heading
\subsection{Research design and evidence hierarchy}\label{sec:2.1}
% END SOURCE BLOCK 036

% BEGIN SOURCE BLOCK 037 paragraph
Matched traffic, fleet and task inputs isolate infrastructure changes under a frozen policy; recorded actions are checked rather than assumed identical. Conclusions concern this evaluator, not retrained or physically validated systems.
% END SOURCE BLOCK 037

% BEGIN SOURCE BLOCK 038 paragraph
The programme is adaptive: E2c excludes its discovery seed and the original morning replication excludes its inspected pilot. Frozen source and accepted records determine experimental meaning; correspondence clarifies intention. Raw audits and kernel explanations are retrospective, distinct from the original protocols \hyperref[source-s0]{[S0]}--\hyperref[source-s8]{[S8]}, \hyperref[source-s13]{[S13]}.
% END SOURCE BLOCK 038

% BEGIN SOURCE BLOCK 039 heading
\subsection{Architecture and information boundaries}\label{sec:2.2}
% END SOURCE BLOCK 039

% BEGIN SOURCE BLOCK 040 paragraph
Figure~\ref{fig:1} separates actor mode choice, environmental radio targeting and infrastructure execution/admission. V2I ingress is the strongest simulated radio link. Ingress execution retains that RSU; load-balancing arms may assign elsewhere. A target remains a proposal until admission \hyperref[source-s4]{[S4]}.
% END SOURCE BLOCK 040

% BEGIN SOURCE BLOCK 041 figure
\begin{figure}[htbp]
\centering
\includesvg[width=\linewidth]{assets/architecture}
\caption{Experimental responsibility boundaries. Only admitted V2I work enters the selected execution queue; logical forwarding applies when ingress and execution differ. Resource scaling is outside the completed comparison.}\label{fig:1}
\end{figure}
% END SOURCE BLOCK 041

% BEGIN SOURCE BLOCK 042 paragraph
The 17-dimensional observation includes task descriptors, vehicle queues, state of charge, radio summaries, compute capability, EV status and observed V2V target tier, but excludes RSU workload. One mode per active vehicle-second serves all independently sampled operational tasks in its five slots. Admission can nevertheless affect transmit energy, then SoC and later observations/actions. Frozen weights therefore do not guarantee fixed actions; the old matched-action finding is empirical, and the new protocol allows this feedback.
% END SOURCE BLOCK 042

% BEGIN SOURCE BLOCK 043 paragraph
V2V selection excludes the source and full-queue peers, then chooses the strongest remaining simulated link. Fading as well as distance determines quality, so this is not nearest-neighbour selection. Peer workload and capability affect latency, not this ranking. Observation and operational links use separate fading samples. Randy confirmed these inherited conventions and one mode/helper per vehicle-second; they were held constant across placement arms \hyperref[source-s4]{[S4]}.
% END SOURCE BLOCK 043

% BEGIN SOURCE BLOCK 044 heading
\subsection{Scenarios, actor and controlled inputs}\label{sec:2.3}
% END SOURCE BLOCK 044

% BEGIN SOURCE BLOCK 045 paragraph
Table~\ref{tab:2} identifies the scenarios. Evaluations import saved positions without rerunning SUMO. Morning input checks reconstruct canonical arrays and verify entry markers because padded slots need not retain one vehicle identity. Queues reset for each new visit; SoC does not. The incident trace retains its earlier convention without the same entry channel \hyperref[source-s7]{[S7]}.
% END SOURCE BLOCK 045

% BEGIN SOURCE BLOCK 046 table
\begingroup
\small
\begin{xltabular}{\textwidth}{@{}YYY@{}}
\caption{Scenario identity and fixed placement controls. Differences between scenarios are retained and documented, not treated as a single-factor intervention.}\label{tab:2} \\
\toprule
Setting & Incident scenario & Morning scenario \\
\midrule\endfirsthead
\toprule
Setting & Incident scenario & Morning scenario \\
\midrule\endhead
\bottomrule\endfoot
Date and local window & 15 March 2024, 20:00--21:00 & 15 October 2024, 08:00--11:00 \\
One-second steps & 3,600 & 10,800 \\
Padded vehicle slots & 2,488 & 215 \\
RSUs & 10 & 9 \\
SUMO seed & 43 & 42 \\
Vehicle queue convention & Conserved, legacy trace convention & Conserved, canonical per-visit reset \\
Placement-study RSU limit & 6,220 tasks per RSU & 6,220 tasks per RSU \\
Service / forwarding / scaling & 1\ensuremath{\times} / 0 ms / off & 1\ensuremath{\times} / 0 ms / off \\
\end{xltabular}
\endgroup
% END SOURCE BLOCK 046

% BEGIN SOURCE BLOCK 047 paragraph
The same archived one-hot-17 checkpoint is used throughout. Its supplied Model C training used synthetic mobility, 128 environments, learning rate 0.003 and seed 100; training was not reproduced here. The provisional UK2030 preset samples RPi/Jetson/GPU tiers with probabilities (0.40, 0.35, 0.25) and EV status with probability 0.22, not measured Manchester shares. Original replications fix evaluator seed 0 while varying fleet seed; the later eight blocks vary both (\S\,\ref{sec:3.4}). Neither varies actor training or trace.
% END SOURCE BLOCK 047

% BEGIN SOURCE BLOCK 048 paragraph
Each active vehicle offers $\min(\operatorname{Poisson}(1.5),5)$ tasks per second; inactive slots offer none. Operational task types are independently sampled with probabilities (0.20, 0.30, 0.50). Types 1--3 have deadlines (100, 500, 100) ms, mean input sizes (1, 0.0012, 0.001) MB and workloads (1,254, 2,100, 15) million cycles. Input sizes multiply their type mean by an independent uniform 0.8--1.2 factor. The five-slot cap changes offered work as well as scheduling opportunities \hyperref[source-s4]{[S4]}.
% END SOURCE BLOCK 048

% BEGIN SOURCE BLOCK 049 paragraph
For homogeneous RSUs, service milliseconds are $s = \frac{\xi \times C_t}{2.2 \times 1.5 \times \min(12,p_t) \times 0.9 \times \min(5,g_t)}$, where $C_t$ is million cycles, $p_t=(4,4,2)$, $g_t=(40,5,1)$, and \ensuremath{\xi} is uniform on 0.9--1.1. The denominator uses 2.2 GHz, 1.5 instructions per cycle, effective cores, utilisation and capped GPU acceleration. Million-cycle/GHz conversion yields milliseconds. Before noise, services are approximately (21.111, 35.354, 2.525) ms. All placement arms use service multiplier 1.
% END SOURCE BLOCK 049

% BEGIN SOURCE BLOCK 050 paragraph
Absolute capacity avoids a confound: the incident ratio 2.5 \ensuremath{\times} 2,488 gives 6,220 tasks per RSU; applying 2.5 to the morning width would give only 538. The morning protocol fixes 6,220 explicitly. Density, duration, layout and vehicle-entry conventions still differ.
% END SOURCE BLOCK 050

% BEGIN SOURCE BLOCK 051 heading
\subsection{Time, queues, admission and outcomes}\label{sec:2.4}
% END SOURCE BLOCK 051

% BEGIN SOURCE BLOCK 052 paragraph
A batch advances simulated time by one second. Five internal task substeps order admissions; they are not five 200 ms clock advances. Work accumulates across them before one service drain. Let $W[r]$ denote remaining service milliseconds at RSU $r$, $Q[r]$ its aggregate task count and C the task limit. These variables serve different purposes: identical counts can represent different work, and increasing C does not accelerate service \hyperref[source-s4]{[S4]}.
% END SOURCE BLOCK 052

% BEGIN SOURCE BLOCK 053 paragraph
For each offered task $i$, let $a[i]$ indicate admission, $y[i]$ indicate the recorded deadline success, and $z[i,c]$ indicate terminal failure category c. The intended accounting contract requires one admission or failure and no success for rejected work. The per-run primary outcome and validation requirements are:
% END SOURCE BLOCK 053

% BEGIN SOURCE BLOCK 054 equation
\begin{equation}\label{eq:1}
 A=\frac{N_{\mathrm{met}}}{N_{\mathrm{offered}}}
  =\frac{\sum_i y[i]}{N_{\mathrm{offered}}};
 \qquad \text{reported percentage}=100A.
\end{equation}
% END SOURCE BLOCK 054

% BEGIN SOURCE BLOCK 055 equation
\begin{equation}\label{eq:2}
 a[i]+\sum_c z[i,c]=1;\qquad y[i]\leq a[i];\qquad
 N_{\mathrm{offered}}=N_{\mathrm{admitted}}+N_{\mathrm{failed}}.
\end{equation}
% END SOURCE BLOCK 055

% BEGIN SOURCE BLOCK 056 paragraph
Here $N_{\mathrm{failed}}$ includes both rejection and unavailability. Categories distinguish local queue rejection, V2V queue rejection or unavailability, and V2I unavailability, backlog-gate rejection or capacity rejection. Gate failure takes precedence over capacity failure when both apply. Selected targets are proposals; rejected tasks neither execute nor reserve work (Figure~\ref{fig:2}). Admitted attainment uses $N_{\mathrm{admitted}}$ instead, so must not replace Equation~\eqref{eq:1}. Offered, admitted and successful-task latency means similarly have different populations.
% END SOURCE BLOCK 056

% BEGIN SOURCE BLOCK 057 figure
\begin{figure}[htbp]
\centering
\includesvg[width=\linewidth]{assets/task_accounting}
\caption{Offered tasks remain in the primary denominator. Admission, execution assignment and deadline outcome are separate fields; selecting an RSU does not establish that work executed there.}\label{fig:2}
\end{figure}
% END SOURCE BLOCK 057

% BEGIN SOURCE BLOCK 058 paragraph
The inherited gate tests existing effective workload, including the candidate offset $O[i]$, against deadline $D[i]$:
% END SOURCE BLOCK 058

% BEGIN SOURCE BLOCK 059 equation
\begin{equation}\label{eq:3}
 W[r]+O[i]<D[i],\quad\text{together with an available task-count place.}
\end{equation}
% END SOURCE BLOCK 059

% BEGIN SOURCE BLOCK 060 paragraph
The inequality is strict. It excludes the candidate's service, radio transfer and forwarding cost, and therefore is a backlog filter rather than an end-to-end feasibility guarantee. Candidate ordering and offsets differ between algorithms (\S\,\ref{sec:2.5}). Admission uses positive ingress link quality as its coarse radio condition; the latency model additionally requires quality at least 0.15 for a usable transmission. Consequently, admission can still lead to a transmission-related deadline miss.
% END SOURCE BLOCK 060

% BEGIN SOURCE BLOCK 061 paragraph
For viable V2I execution, the modelled response latency is:
% END SOURCE BLOCK 061

% BEGIN SOURCE BLOCK 062 equation
\begin{equation}\label{eq:4}
 \begin{aligned}
 L_{\mathrm{raw}}[i]={}&T(\mathrm{input}[i],q,c)+W[r]+O[i]+s[i,r]\\
 &+T(0.001,q,c)+f\cdot\mathbf{1}(r\neq\mathrm{ingress}[i]).
 \end{aligned}
\end{equation}
% END SOURCE BLOCK 062

% BEGIN SOURCE BLOCK 063 paragraph
$T(b,q,c)$ is $0.1 + 1{,}000 \times 8b/c$ milliseconds for $q\geq 0.15$ and capacity $c>0.000000001$ Mbps, with $b$ in MB; otherwise it uses a $10^{9}$ ms failure value (up to the 0.1 ms propagation addition). The 0.001 MB return term reuses the ingress quality and capacity. It is not a separately simulated downlink or inter-RSU route. The fixed forwarding charge f is added once. Local latency is vehicle backlog plus local service; V2V adds input/return transfers on the selected peer link, peer backlog, preceding local/V2V reservations and peer service. Service $s$ derives from task cycles, device frequency, parallelism, utilisation and acceleration, with a sampled 0.9--1.1 multiplier. The placement code reserves the same sampled work used for latency: exact service knowledge is a declared information assumption.
% END SOURCE BLOCK 063

% BEGIN SOURCE BLOCK 064 paragraph
The gate-enabled paths pass final admission into latency scoring. Historical \texttt{off} instead passes earlier coarse eligibility while enqueueing uses refined admission: an in-substep capacity rejection need not receive penalty latency. Its historical numerical impact remains unquantified (\S\,\ref{sec:3.1}). The final mapping is $L[i]=10D[i]$ when $L_{\mathrm{raw}}[i]\geq 10^{8}$ ms, and $L[i]=L_{\mathrm{raw}}[i]$ otherwise. Finite late latencies are not generally clipped; success uses $L[i]\leq D[i]$, inclusively. Finite rejection latency alone is not a false deadline success.
% END SOURCE BLOCK 064

% BEGIN SOURCE BLOCK 065 paragraph
For each queue over an interval, work accounting requires:
% END SOURCE BLOCK 065

% BEGIN SOURCE BLOCK 066 equation
\begin{equation}\label{eq:5}
 W_{\mathrm{initial}}+\sum\text{admitted service}
 =\sum\text{served service}+W_{\mathrm{final}}.
\end{equation}
% END SOURCE BLOCK 066

% BEGIN SOURCE BLOCK 067 paragraph
After the fifth substep, $S[r]=\min(W_{\mathrm{post}}[r],1{,}000)$ is served and $W_{\mathrm{next}}[r]=W_{\mathrm{post}}[r]-S[r]$. If this empties the queue, its task count becomes zero. Otherwise the evaluator subtracts $\left\lfloor Q_{\mathrm{post}}[r]\times S[r]/W_{\mathrm{post}}[r]\right\rfloor$, retaining at least one task while positive work remains. This aggregate carry approximation has no individual departure-event model. Equation~\eqref{eq:5} checks the declared dynamics, not their physical calibration.
% END SOURCE BLOCK 067

% BEGIN SOURCE BLOCK 068 heading
\subsection{Placement algorithms and their implementation}\label{sec:2.5}
% END SOURCE BLOCK 068

% BEGIN SOURCE BLOCK 069 paragraph
Table~\ref{tab:3} separates execution destination from the gate switch. The identifier \texttt{jsq} actually selects least service workload, not shortest task count. \texttt{dla} combines this common-target selector with Equation~\eqref{eq:3}; these inherited names do not define canonical scheduling algorithms \hyperref[source-s2]{[S2]}, \hyperref[source-s4]{[S4]}.
% END SOURCE BLOCK 069

% BEGIN SOURCE BLOCK 070 table
\begingroup
\small
\begin{xltabular}{\textwidth}{@{}YYYY@{}}
\caption{Infrastructure conditions. All retain strongest-link radio ingress. ``Live'' eligibility is refreshed at each substep; the historical off mode retains step-entry coarse eligibility.}\label{tab:3} \\
\toprule
Mode & Execution destination & Backlog gate & Coarse eligibility \\
\midrule\endfirsthead
\toprule
Mode & Execution destination & Backlog gate & Coarse eligibility \\
\midrule\endhead
\bottomrule\endfoot
\texttt{off} & Ingress & Off & Step entry \\
\texttt{jsq} & Common least-workload target & Off & Live \\
\texttt{ingress\_dla} & Ingress & On & Live \\
\texttt{dla} & Common least-workload target & On & Live \\
\texttt{per\_task\_dla} & Causal least-workload target & On & Live \\
\texttt{causal\_round\_robin} & Persistent cyclic target & On & Live \\
\end{xltabular}
\endgroup
% END SOURCE BLOCK 070

% BEGIN SOURCE BLOCK 071 paragraph
Both placement algorithms below operate once per task substep, using its entry state $W$, Q. Index $i$ follows ascending padded-vehicle order, not arrival-time or deadline priority. Each candidate has actual service $s[i,r]$, deadline $D[i]$ and a fixed actor action and radio ingress. Every argmin considers all RSUs and breaks ties by lowest index; a full selected RSU is rejected rather than excluded and retried elsewhere.
% END SOURCE BLOCK 071

% BEGIN SOURCE BLOCK 072 paragraph
For common-target reconciliation, let $E[i]$ mean active V2I attempt, positive ingress quality and $Q[r[i]]<C$ at substep entry. For an admission mask $M$, define $O[i,M]$ as the sum of $s[j,r[i]]$ over earlier $j<i$ with the same target and $M[j]$ true; $H[i,M]$ is their count. Offsets are exclusive: a candidate's own service is omitted. Algorithm~\ref{alg:1} performs exactly three simultaneous replacements. Proposition~\ref{prop:1} later supplies an exact-model guarantee under stated assumptions; it is not a blanket float32 guarantee.
% END SOURCE BLOCK 072

% BEGIN SOURCE BLOCK 073 algorithm
\begin{algorithm}[htbp]
\caption{common-target placement with gate (one substep)}\label{alg:1}
\begin{Verbatim}[fontsize=\small,breaklines=true]
r[i] = lowest-index argmin(W), shared by all candidates
M = E
repeat exactly three times:
    compute O[i,M] and H[i,M] from the previous whole mask
    M_new[i] = E[i] and (W[r[i]] + O[i,M] < D[i])
                       and (H[i,M] < C - Q[r[i]])
    M = M_new simultaneously
recompute final offsets O[i,M] for latency
label E-candidates failing the final workload test as gate failures
label remaining E-candidates excluded by M as capacity failures
commit each M-admitted task and its full service once
unavailable or rejected tasks reserve nothing
\end{Verbatim}
\end{algorithm}
% END SOURCE BLOCK 073

% BEGIN SOURCE BLOCK 074 paragraph
Target selection never changes during reconciliation. The retained mask and recomputed offsets determine recorded processing. Section~\ref{sec:3.5} distinguishes the exact guarantee from its float32 boundary. Ingress with the gate uses the same offset/reconciliation machinery with individual ingress targets instead of the common argmin. Gate-off modes omit the workload test.
% END SOURCE BLOCK 074

% BEGIN SOURCE BLOCK 075 algorithm
\begin{algorithm}[htbp]
\caption{causal per-task placement with gate (one substep)}\label{alg:2}
\begin{Verbatim}[fontsize=\small,breaklines=true]
repeat exactly three complete passes:
    working_W = W; working_Q = Q       # restart, not cumulative
    for i in ascending padded-vehicle order:
        r[i] = lowest-index argmin(working_W)
        O[i] = working_W[r[i]] - W[r[i]]
        require active V2I, positive quality and Q[r[i]] < C
        test working_W[r[i]] < D[i], then working_Q[r[i]] < C
        if admitted:
            working_W[r[i]] += s[i,r[i]]
            working_Q[r[i]] += 1
        otherwise label failure; reserve nothing
retain the last pass; commit its admitted work once
\end{Verbatim}
\end{algorithm}
% END SOURCE BLOCK 075

% BEGIN SOURCE BLOCK 076 paragraph
The causal gate and capacity see earlier admissions immediately. Repeated passes restart identically and commit once, before the next substep; neither triples reservations nor drains service between substeps. The comparison therefore changes selection frequency, reservation visibility and reconciliation semantics, not argmin frequency alone.
% END SOURCE BLOCK 076

% BEGIN SOURCE BLOCK 077 paragraph
For an illustrative scheduling calculation, assume three empty identical RSUs, four viable 40 ms tasks with 100 ms deadlines, nonbinding capacity and negligible communication. Common-target chooses RSU 0; its masks settle at three admissions and one gate rejection. The admitted offsets are 0, 40 and 80 ms, leaving 120, 0, 0 ms: the third admitted task misses its deadline despite passing the backlog gate. Per-task selects 0, 1, 2, 0, leaves 80, 40, 40 ms and all four meet deadlines under these assumptions. Figure~\ref{fig:3} depicts the proposals. This constructed example explains mechanics; it is neither a recorded task group nor an estimate of recoverable historical failures.
% END SOURCE BLOCK 077

% BEGIN SOURCE BLOCK 078 figure
\begin{figure}[htbp]
\centering
\includesvg[width=\linewidth]{assets/dispatch_example}
\caption{Four illustrative proposals. The common destination remains fixed through reconciliation; causal placement updates after admissions. The accompanying worked example states admissions and deadline outcomes under deliberately simplified assumptions.}\label{fig:3}
\end{figure}
% END SOURCE BLOCK 078

% BEGIN SOURCE BLOCK 079 heading
\subsection{Experimental progression and inference}\label{sec:2.6}
% END SOURCE BLOCK 079

% BEGIN SOURCE BLOCK 080 paragraph
E0 establishes the recorded measurement foundation. E1 compares three queue limits at fixed service over five matched draws; E2/E2b explore placement and admission on one draw. E2c compares common-target with ingress on new incident seeds 1--4. After inspecting them, E2d adds per-task placement on those same draws and reuses their controls: an adaptive extension, not new independent observations \hyperref[source-s0]{[S0]}--\hyperref[source-s3]{[S3]}.
% END SOURCE BLOCK 080

% BEGIN SOURCE BLOCK 081 paragraph
The morning pilot uses seed 1 for ingress and per-task placement. After inspecting it, the replication declares seeds 0, 2, 3 and 4 and all three arms, yielding twelve new full runs. The primary contrasts are per-task minus ingress, common-target minus ingress, and per-task minus common-target. The pilot is excluded; a supplementary two-arm five-draw summary is descriptive. The local protocol and manifest preceded the new outcomes, but there was no external preregistration \hyperref[source-s7]{[S7]}, \hyperref[source-s8]{[S8]}.
% END SOURCE BLOCK 081

% BEGIN SOURCE BLOCK 082 paragraph
For each contrast, a fleet draw supplies one paired difference in percentage points. Draws receive equal weight. The mean difference has a Student-$t$ interval, $\text{mean}\pm t\times\frac{\text{sample standard deviation}}{\sqrt{n}}$. Incident E2c/E2d retain their reported individual 95\% intervals. The morning protocol additionally uses Bonferroni simultaneous intervals for three contrasts, with critical value $t$ at probability $1-0.05/(2\times 3)$, three degrees of freedom. Its reversal criterion requires the simultaneous common-target-minus-ingress interval below zero and the per-task-minus-ingress interval above zero. Four draws cannot strongly diagnose distributional assumptions; these small-sample intervals remain conditional parametric summaries, not task-level significance tests.
% END SOURCE BLOCK 082

% BEGIN SOURCE BLOCK 083 heading
\subsection{Sensitivities, mechanism audit and validation}\label{sec:2.7}
% END SOURCE BLOCK 083

% BEGIN SOURCE BLOCK 084 paragraph
The state-information pilot uses incident seed 1 and report ages 0, 100, 500 and 1,000 ms, alongside ingress. Admission remains live. For previous post-admission workload B and positive age $d$, the report is $\max(B-(1{,}000-d),0)$, with current-batch reservations added immediately. This assumes continuous service between batches without redistributing arrivals. Startup uses live state and recorded age zero \hyperref[source-s5]{[S5]}.
% END SOURCE BLOCK 084

% BEGIN SOURCE BLOCK 085 paragraph
Four ten-step forwarding probes qualify fixed-latency transformations of saved full records: charges leave admission, queues and actions unchanged. No full positive-cost simulation or network model is implied; penalty-equal ambiguity remains unresolved \hyperref[source-s6]{[S6]}.
% END SOURCE BLOCK 085

% BEGIN SOURCE BLOCK 086 paragraph
The completed September audit establishes admission independently of success: active masks define offers, enqueue fields identify V2I admission, and non-V2I rejection predicates identify terminal failures. It cross-checks flags, categories, latency/deadlines, assignments and summaries. Execution fields describe admitted assignment, not individual departures \hyperref[source-s13]{[S13]}.
% END SOURCE BLOCK 086

% BEGIN SOURCE BLOCK 087 paragraph
Task joins require matched trace-row/substep/slot coordinates and corresponding fleet, task, arrival, action and ingress fields. Unsaved sizes correspond through the arm-independent key schedule. Gains minus losses must equal the success difference over the common offered denominator; later queues can diverge, so these are descriptive transitions.
% END SOURCE BLOCK 087

% BEGIN SOURCE BLOCK 088 paragraph
Retrospective diagnostics distinguish simultaneous common-target admission A, causal admission replaying A's targets B, and causal reselection C. The completed proof and numerical investigation are retained in Appendix~\ref{app:C}, separately from observed fleet results \hyperref[source-s14]{[S14]}.
% END SOURCE BLOCK 088

% BEGIN SOURCE BLOCK 089 heading
\subsection{Cyclic comparator and retrospective measurement design}\label{sec:2.8}
% END SOURCE BLOCK 089

% BEGIN SOURCE BLOCK 090 paragraph
Causal round-robin isolates workload-aware selection from simple spreading. A pointer starts at RSU 0 and persists across candidates, substeps and batches. Each active V2I attempt with positive ingress quality proposes the pointer's RSU and advances it modulo $R$, even if capacity or the strict gate subsequently rejects it. Unavailable radio and padding do not advance it; there is no retry or full-node skipping. Candidate order, sampled work, gate/capacity rules, precision, accounting, forwarding and drain match causal per-task placement. Repeated passes restart from the same pointer and queues; the final pointer commits once. A versioned evaluator copy preserves the frozen implementations. Qualification checked unchanged-arm compatibility, shared inputs, final-admission accounting and pointer restart; Appendix~\ref{app:E} records its finite coverage \hyperref[source-s15]{[S15]}, \hyperref[source-s16]{[S16]}.
% END SOURCE BLOCK 090

% BEGIN SOURCE BLOCK 091 paragraph
Scheduler timing uses actual JAX helpers, including three-replacement prefixes; precomputed candidate work replaces only the service provider. Sixteen configurations combine four arms, $N/R=215/9$ or 2,488/10, $K=5$ and sparse/empty or busy/nonempty fixtures. Arms share arrays with coupled type/deadline/service values. Busy entry queues are imposed stress states, excluded by exact clearing from empty initialisation (\S\,\ref{sec:3.5}). Following JAX guidance, compilation is separate, inputs are device-resident, five warm-ups precede 100 synchronised calls. Actor, radio, service generation, transfer, scoring and I/O are excluded. Repetitions measure host variability, not fleet uncertainty \cite{ref14}.
% END SOURCE BLOCK 091

% BEGIN SOURCE BLOCK 092 paragraph
Retrospective type aggregation uses the eight authenticated ingress/per-task files from the original four-draw morning sample. Categories supply admission, gate rejection and admitted misses; type totals must reproduce all-task summaries and paired net changes. All types and draws are retained, without subgroup significance tests.
% END SOURCE BLOCK 092

% BEGIN SOURCE BLOCK 093 heading
\section{Evaluation and Reflection}\label{sec:3}
% END SOURCE BLOCK 093

% BEGIN SOURCE BLOCK 094 heading
\subsection{Accounting validity and the capacity question}\label{sec:3.1}
% END SOURCE BLOCK 094

% BEGIN SOURCE BLOCK 095 paragraph
E0 addressed a concrete mismatch between task scoring and queue admission. In the legacy clamp path, the evaluator scored eligible V2I tasks and multiplied their combined service work by the fraction fitting the task-count ceiling. Surplus tasks were not individually identified in earlier scoring. The 5 August repair introduced per-task admission/rejection and full enqueueing of admitted work. Source history credits Randy's implementation and Abdulla's motivating accounting analysis; E0 qualified recorded accounting in that configuration \hyperref[source-s0]{[S0]}, \hyperref[source-s12]{[S12]}.
% END SOURCE BLOCK 095

% BEGIN SOURCE BLOCK 096 paragraph
Illustratively, suppose one queue place remains and two viable tasks each require 40 ms. Legacy fractional enqueue adds (1/2) \ensuremath{\times} 80 = 40 ms, although both can remain in the scored population. Reject-mode admission instead accepts one task in order, enqueues its full 40 ms and assigns the other a failure. With unequal services, fractional enqueue can also differ from the work of the particular admitted task. This is a constructed explanation, not an observed historical pair.
% END SOURCE BLOCK 096

% BEGIN SOURCE BLOCK 097 paragraph
The relevant invariant is that offers partition into admissions and terminal failures, with no execution or success assigned to rejected work. E0's negative fixture changes the summary offered count without adding a task record and must fail. Summary work partition alone is weaker: rejected work is calculated as offered minus admitted. Later reconstruction instead checks enqueue, drain and queue endpoints independently.
% END SOURCE BLOCK 097

% BEGIN SOURCE BLOCK 098 paragraph
Archived E0 receipts report passing their checks, but source inspection identifies a remaining exception in E0/E1 and E2's reused off path: coarse eligibility reaches latency scoring while refined admission governs enqueueing. Recorded capacity rejections establish exposure, not false successes. Original E0/E1/E2 outputs are unavailable following author-reported deletion, with no known backup. Their numerical scoring impact is not assessable from surviving summaries and receipts; Appendix~\ref{app:B} distinguishes this gap from later audit coverage \hyperref[source-s13]{[S13]}.
% END SOURCE BLOCK 098

% BEGIN SOURCE BLOCK 099 paragraph
E1 varied waiting-room capacity across five matched draws. The archived score difference for 99,520 versus 1,866 tasks per RSU was \ensuremath{-}0.01094 percentage points, with a 95\% interval [\ensuremath{-}0.02465, +0.00276]. Recalculation from surviving run summaries reproduces this interval; it is not fresh task-level validation. The archived E1 score contrast is statistically inconclusive under the historical implementation. Its interpretation as admission-consistent deadline attainment remains qualified by the unquantified scoring-mask exception. The interval describes sampling uncertainty conditional on that implementation, not the additional measurement uncertainty. Neither equivalence nor statistically supported degradation follows \hyperref[source-s1]{[S1]}.
% END SOURCE BLOCK 099

% BEGIN SOURCE BLOCK 100 paragraph
The saved secondary summaries report more admissions and lower rejection at larger limits, while mean admitted latency rises from approximately 3.77 to 12.19 and 132.49 seconds. These changing admitted populations cannot substitute for the offered-task contrast; their latency summaries retain the scoring-mask qualification. More retained work need not mean more useful service.
% END SOURCE BLOCK 100

% BEGIN SOURCE BLOCK 101 paragraph
Changing the off scoring mask could affect transmit energy, SoC and later observations: even admission-consistent saved flags would describe fixed history, not a corrected rerun. Separately, the completed audit found no non-admitted successes or non-penalty rejection latencies in 19 authenticated September gate-enabled runs, with admission/category/summary agreement. It supports those outcomes without resolving historical measurement impact.
% END SOURCE BLOCK 101

% BEGIN SOURCE BLOCK 102 heading
\subsection{From admission decomposition to the incident reversal}\label{sec:3.2}
% END SOURCE BLOCK 102

% BEGIN SOURCE BLOCK 103 paragraph
E2's exploratory seed-0 comparison showed offered attainment of approximately 68.362\% for strongest-link execution without the deadline gate, 67.568\% for inherited least-busy placement without the gate, and 69.494\% for inherited least-busy placement with the gate. The last comparison changed both placement and admission relative to the strongest-link reference. Its improvement could not establish that load balancing itself helped \hyperref[source-s2]{[S2]}.
% END SOURCE BLOCK 103

% BEGIN SOURCE BLOCK 104 paragraph
E2b added ingress execution with the same deadline gate, reaching approximately 71.577\%. The strongest-link contrast was +3.215 points and the common-target contrast +1.926 points, giving an interaction of \ensuremath{-}1.290 points. However, \texttt{off} uses step-entry coarse saturation eligibility while \texttt{ingress\_dla} refreshes eligibility per substep. Their contrast and interaction bundle the gate with this timing difference and retain the legacy scoring-mask exception (\S\,\ref{sec:2.4}). The manifest records 138 unavailable V2I attempts among 13,076,234 offered tasks in off; that count does not bound downstream effects. E2b identifies a stronger comparator without assigning its entire gain solely to the gate. Subsequent gate-held placement comparisons use live eligibility and refined scoring masks in all arms. In the exploratory draw, common-target remained 2.083 points below ingress with the gate.
% END SOURCE BLOCK 104

% BEGIN SOURCE BLOCK 105 paragraph
E2c tested the common-target comparison on four new incident fleet draws, excluding discovery seed 0. Its mean difference from ingress was \ensuremath{-}2.122 points, with an individual 95\% interval of [\ensuremath{-}2.234, \ensuremath{-}2.011]. All four paired differences were negative. Inspection then established that the inherited scheduler chose one common target per task substep. E2d introduced the causal per-task implementation under the same actor, trace, gate predicate, queue ceiling, service and forwarding controls. Table~\ref{tab:4} retains all four draw-level results \hyperref[source-s3]{[S3]}.
% END SOURCE BLOCK 105

% BEGIN SOURCE BLOCK 106 table
\begingroup
\small
\begin{xltabular}{\textwidth}{@{}YYYYY@{}}
\caption{Incident offered-task deadline attainment, expressed as percentages. E2d reused the exact E2c ingress and common-target controls; these are four paired draws, not separate sets of independent controls.}\label{tab:4} \\
\toprule
Fleet seed & Ingress & Common-target & Per-task & Per-task \ensuremath{-} ingress, pp \\
\midrule\endfirsthead
\toprule
Fleet seed & Ingress & Common-target & Per-task & Per-task \ensuremath{-} ingress, pp \\
\midrule\endhead
\bottomrule\endfoot
1 & 72.003 & 69.794 & 72.467 & +0.464 \\
2 & 70.369 & 68.317 & 70.956 & +0.587 \\
3 & 70.298 & 68.153 & 70.805 & +0.507 \\
4 & 70.887 & 68.805 & 71.438 & +0.551 \\
\end{xltabular}
\endgroup
% END SOURCE BLOCK 106

% BEGIN SOURCE BLOCK 107 paragraph
Per-task exceeds ingress by 0.527 points, individual 95\% interval [+0.442, +0.612], and common-target by 2.649 points, interval [+2.621, +2.678]. RQ2 therefore yields opposite rankings for the two implementations. These intervals summarise an adaptive extension on examined draws; the morning protocol supplies prospective replication.
% END SOURCE BLOCK 107

% BEGIN SOURCE BLOCK 108 paragraph
Ingress determines the practical scale: against 70.889\% mean attainment, +0.527 points means about 53 extra successes per 10,000 offered, or a 0.744\% relative increase. The larger +2.649-point common-target contrast diagnoses that implementation; it is not the gain over the stronger reference. These are descriptive conversions of the same mean, and neither arm establishes optimal dispatch.
% END SOURCE BLOCK 108

% BEGIN SOURCE BLOCK 109 heading
\subsection{Morning replication with the inspected pilot excluded}\label{sec:3.3}
% END SOURCE BLOCK 109

% BEGIN SOURCE BLOCK 110 paragraph
The inspected morning pilot offered 1,744,116 tasks and produced 91.075\% ingress versus 94.028\% per-task attainment: +2.953 points or 51,503 successes. It motivated replication but remains exploratory \hyperref[source-s7]{[S7]}.
% END SOURCE BLOCK 110

% BEGIN SOURCE BLOCK 111 paragraph
The morning protocol declared four new seeds and all three arms before outcomes. Each draw reproduces common-target < ingress < per-task. Mean attainment is 85.655\%, 88.887\% and 92.785\%, respectively; its lower ingress/per-task means than the pilot underline the need for replication \hyperref[source-s8]{[S8]}.
% END SOURCE BLOCK 111

% BEGIN SOURCE BLOCK 112 table
\begingroup
\small
\begin{xltabular}{\textwidth}{@{}YYYYY@{}}
\caption{Morning primary replication. Seed 1 is absent because it was the inspected pilot; each row is a matched fleet draw on the same morning trace.}\label{tab:5} \\
\toprule
Fleet seed & Ingress, \% & Common-target, \% & Per-task, \% & Per-task \ensuremath{-} ingress, pp \\
\midrule\endfirsthead
\toprule
Fleet seed & Ingress, \% & Common-target, \% & Per-task, \% & Per-task \ensuremath{-} ingress, pp \\
\midrule\endhead
\bottomrule\endfoot
0 & 88.817 & 85.660 & 92.705 & +3.888 \\
2 & 87.536 & 83.537 & 92.060 & +4.524 \\
3 & 89.082 & 85.862 & 92.978 & +3.896 \\
4 & 90.113 & 87.559 & 93.399 & +3.286 \\
\end{xltabular}
\endgroup
% END SOURCE BLOCK 112

% BEGIN SOURCE BLOCK 113 paragraph
Table~\ref{tab:6}'s simultaneous per-task-minus-ingress interval is positive and common-target-minus-ingress interval negative, satisfying the prespecified reversal criterion. Figure~\ref{fig:4} displays these primary results.
% END SOURCE BLOCK 113

% BEGIN SOURCE BLOCK 114 table
\begingroup
\small
\begin{xltabular}{\textwidth}{@{}YYYY@{}}
\caption{Morning paired contrasts in percentage points. Individual intervals are two-sided 95\% Student-$t$ intervals; simultaneous intervals use the prespecified three-contrast Bonferroni adjustment. Replication uses four fleet draws.}\label{tab:6} \\
\toprule
Contrast & Mean & Individual 95\% interval & Simultaneous 95\% interval \\
\midrule\endfirsthead
\toprule
Contrast & Mean & Individual 95\% interval & Simultaneous 95\% interval \\
\midrule\endhead
\bottomrule\endfoot
Per-task \ensuremath{-} ingress & +3.899 & [+3.094, +4.703] & [+2.671, +5.126] \\
Common-target \ensuremath{-} ingress & \ensuremath{-}3.232 & [\ensuremath{-}4.176, \ensuremath{-}2.289] & [\ensuremath{-}4.672, \ensuremath{-}1.793] \\
Per-task \ensuremath{-} common-target & +7.131 & [+5.385, +8.877] & [+4.466, +9.795] \\
\end{xltabular}
\endgroup
% END SOURCE BLOCK 114

% BEGIN SOURCE BLOCK 115 figure
\begin{figure}[htbp]
\centering
\includesvg[width=\linewidth]{assets/primary_replication}
\caption{Redrawn from archived morning primary data. Left: within-draw attainment. Right: draw differences, mean diamonds and three-contrast Bonferroni intervals. The inspected pilot is excluded; uncertainty concerns four newly declared fleet draws, conditional on the frozen actor and morning scenario \hyperref[source-s8]{[S8]}.}\label{fig:4}
\end{figure}
% END SOURCE BLOCK 115

% BEGIN SOURCE BLOCK 116 paragraph
The supplementary five-draw ingress/per-task mean, including the pilot, is +3.709 points descriptively; it does not enter Table~\ref{tab:6}. No common-target pilot result is available.
% END SOURCE BLOCK 116

% BEGIN SOURCE BLOCK 117 paragraph
The morning gain is about 390 extra successes per 10,000 offered, a 4.386\% relative increase over ingress. Its larger magnitude cannot be attributed solely to density: RSU layout/count, date, duration, padded fleet width and entry conventions also change. Separate analyses support recurrence in two Manchester scenarios, without estimating a pooled geographical effect.
% END SOURCE BLOCK 117

% BEGIN SOURCE BLOCK 118 heading
\subsection{Confirmation under joint randomness and comparison with cyclic spreading}\label{sec:3.4}
% END SOURCE BLOCK 118

% BEGIN SOURCE BLOCK 119 paragraph
All 32 full morning cells and eight block-control receipts passed after bounded qualification; no attempt failed or was retried. Fleet/evaluator pairs (100,200) through (107,207), controls, rotating arm order and analysis were sealed before full outcomes. Block 0 remains included; the original sample and pilots are not pooled \hyperref[source-s16]{[S16]}.
% END SOURCE BLOCK 119

% BEGIN SOURCE BLOCK 120 paragraph
Evaluator seed varies observation descriptors, arrivals, operational types/sizes, fading and service wobble; fleet seed varies tier, EV status and initial SoC. All four arms shared the required exogenous arrays exactly within each block. Actions matched exactly in all 24 non-reference arm/block comparisons. Small observation and logit differences occurred in block 1; they did not change actions. The estimand remains the total scheduling intervention under frozen weights. Observed action matching extends the original finding to these streams without forcing replay or guaranteeing future equality.
% END SOURCE BLOCK 120

% BEGIN SOURCE BLOCK 121 table
\begingroup
\small
\begin{xltabular}{\textwidth}{@{}YYY@{}}
\caption{New morning confirmation, eight paired joint-seed blocks with equal weighting. Two-sided Bonferroni simultaneous 95\% Student-$t$ intervals cover the three predeclared contrasts; positive differences favour the first arm.}\label{tab:7} \\
\toprule
Contrast & Mean difference, pp & Simultaneous 95\% interval, pp \\
\midrule\endfirsthead
\toprule
Contrast & Mean difference, pp & Simultaneous 95\% interval, pp \\
\midrule\endhead
\bottomrule\endfoot
Per-task \ensuremath{-} ingress & +4.137 & [+3.556, +4.717] \\
Common-target \ensuremath{-} ingress & \ensuremath{-}3.504 & [\ensuremath{-}3.980, \ensuremath{-}3.028] \\
Per-task \ensuremath{-} round-robin & +0.631 & [+0.511, +0.750] \\
\end{xltabular}
\endgroup
% END SOURCE BLOCK 121

% BEGIN SOURCE BLOCK 122 paragraph
The declared reversal criterion holds, evaluated using both directional requirements within the same simultaneous family. The mean per-task difference corresponds to +413.7 successes per 10,000 offered tasks relative to ingress. Workload-aware targeting outperformed cyclic spreading in this study. Its mean benefit is 63.1 successes per 10,000 offered tasks over round-robin, accompanying higher scheduler-only cost on the separate fixtures in Table~\ref{tab:10}.
% END SOURCE BLOCK 122

% BEGIN SOURCE BLOCK 123 paragraph
The critical value is 3.127552, recomputed with seven degrees of freedom. Independent joint blocks and approximately normal paired effects remain assumptions; eight blocks cannot strongly diagnose tails. Joint randomness extends the original fleet-conditional result within this morning scenario. It does not identify a unique mechanism decomposition or generalise across geography, actor training, physical timing or distributed information. Full counts, control diagnostics and elapsed times are in Appendix~\ref{app:E}.
% END SOURCE BLOCK 123

% BEGIN SOURCE BLOCK 124 heading
\subsection{Mechanism explanation and task-outcome interpretation}\label{sec:3.5}
% END SOURCE BLOCK 124

% BEGIN SOURCE BLOCK 125 paragraph
Original morning common-target work used only RSUs 0--4, with mean shares approximately 48.34\%, 30.33\%, 14.63\%, 5.21\% and 1.48\%; RSUs 5--8 received none. Per-task shares were close to one ninth each. Figure~\ref{fig:5} describes allocation, without establishing that equal workloads caused the attainment difference \hyperref[source-s8]{[S8]}.
% END SOURCE BLOCK 125

% BEGIN SOURCE BLOCK 126 figure
\begin{figure}[htbp]
\centering
\includesvg[width=\linewidth]{assets/rsu_workload_distribution}
\caption{Admitted service-work distribution, redrawn from archived CSVs. Bars are means; whiskers are observed minima/maxima over four fleet draws, not confidence intervals. Aggregate shares motivated a separate post-hoc audit; they do not by themselves establish intermediate scheduling decisions or counterfactual task outcomes.}\label{fig:5}
\end{figure}
% END SOURCE BLOCK 126

% BEGIN SOURCE BLOCK 127 paragraph
The post-hoc audit found zero saved end-of-second workload and task counts in all four common-target runs, and target k at substep k whenever a V2I attempt exposed it. Empty substeps do not reveal the internal selector. Intermediate workloads were reconstructed, not saved directly, using recorded admissions/destinations and the qualified service reference. All endpoints and observable target minima agreed, although zero endpoints alone weakly constrain intermediate paths. Reconstructed pre-drain workload stayed below 538.730 ms and task counts below 26, far from the 1,000 ms drain and 6,220-task ceiling \hyperref[source-s9]{[S9]}.
% END SOURCE BLOCK 127

% BEGIN SOURCE BLOCK 128 paragraph
A conditional deduction from Algorithm~\ref{alg:1} explains the concentration. With $K$ substeps and $R$ RSUs, at most $\min(K,R)$ distinct destinations receive new assignments in one batch: each substep contributes at most one target, so their union has at most $K$ members. This bound does not require empty queues or a particular tie rule, and concerns new assignments rather than all work already present.
% END SOURCE BLOCK 128

% BEGIN SOURCE BLOCK 129 paragraph
A recurring low-index prefix additionally requires empty batch starts, positive productive admissions, lowest-index ties and no intervening drain. RSU 0 is selected first; its new work makes the next empty index preferable. Induction gives a prefix when $R\geq K$, while an unproductive substep does not advance it. End-of-batch clearing resets the tie. Without clearing or deterministic ties, the per-batch bound could hold while the full-run union covered every RSU.
% END SOURCE BLOCK 129

% BEGIN SOURCE BLOCK 130 paragraph
The arrays recorded 511,684 deadline-gate rejections and zero RSU-capacity rejections. In reconstructed states, every gate rejection coexisted with at least five idle RSUs possessing spare task capacity. Larger waiting rooms would not address that gate bottleneck. Spare capacity does not establish that each task would succeed elsewhere: redistribution changes later queues, and end-to-end latency also matters.
% END SOURCE BLOCK 130

% BEGIN SOURCE BLOCK 131 proposition
\begin{proposition}[fixed-target admission in the studied exact model]\label{prop:1}
Fix candidate order, eligibility, destinations, initial nonnegative workload/count and nonnegative service demands, with strict backlog and exclusive count-capacity tests and no intervening drain. At each nonempty destination, let $n$ be the eligible candidate count and $d$ the number of distinct deadlines. Starting from the eligibility mask, simultaneous replacement reaches the unique causal fixed point within $\min(n,2d-1)$ replacements; destinations decouple.\end{proposition}
% END SOURCE BLOCK 131

% BEGIN SOURCE BLOCK 132 proof
\begin{proof}[Proof sketch]
Earlier-candidate dependence yields the unique causal solution and the candidate-count bound. The update reverses inclusion. Removing the first excess admission either exhausts capacity or fixes rejection of all later deadlines at or below its threshold. Every two replacements therefore eliminate at least one unresolved deadline class; one class needs one replacement. Appendix~\ref{app:C} gives the full induction and assumptions \hyperref[source-s14]{[S14]}.\end{proof}
% END SOURCE BLOCK 132

% BEGIN SOURCE BLOCK 133 paragraph
Here ``production'' means the studied evaluator configuration, not deployed roadside hardware. Its two deadline values imply three replacements suffice in exact arithmetic, including nonempty queues and capacity limits. The earlier three-threshold instability is outside this domain. This rules out unrestricted reconciliation failure as an exact-model explanation, while leaving target reselection and executable rounding distinct.
% END SOURCE BLOCK 133

% BEGIN SOURCE BLOCK 134 paragraph
The same model makes clearing structural from empty initialisation. The last admitted task sees less than 500 ms of backlog and contributes at most approximately 38.889 ms of service. Each queue therefore remains below 538.889 ms before a 1,000 ms drain. The observed zero endpoints and reconstructed maxima near 538 ms agree with this consequence. Queue clearing and the resulting destination prefix thus follow from the stated gate/service/timing assumptions, not traffic density alone. This deduction is conditional on exact admission arithmetic and service multiplier 1.
% END SOURCE BLOCK 134

% BEGIN SOURCE BLOCK 135 paragraph
Float32 limits executable equivalence. Inclusive prefix sums followed by subtraction can round differently from causal accumulation, changing a strict gate even for sub-0.001 ms differences. The completed numerical investigation includes discrete disagreements and a constructed alternating mask; that task misses its deadline under either treatment. These fixtures neither establish full-trajectory reachability nor measure historical prevalence. Appendix~\ref{app:C} retains the intermediate masks, threshold values and unsuccessful transfers, alongside the exact proof.
% END SOURCE BLOCK 135

% BEGIN SOURCE BLOCK 136 paragraph
Reselection can also lose completions. In the two-RSU illustration, both arms admit eight tasks; fixed-target admission completes eight and per-task reselection seven. Spreading earlier long-deadline work makes a later 100 ms task finish at 112.929 rather than 91.818 ms. Current workload omits future deadlines and the gate omits own service. Table~\ref{tab:C1} preserves the calculation: production task values, reduced infrastructure, zero transfers, and B replaying A targets rather than acting autonomously. Twelve transfers to the studied dimensions retained no loss. This identifies a possible interaction, not the cause or frequency of recorded losses; proposed full-evaluator prefixes remain unrun \hyperref[source-s10]{[S10]}, \hyperref[source-s14]{[S14]}.
% END SOURCE BLOCK 136

% BEGIN SOURCE BLOCK 137 paragraph
\textbf{Task outcomes in the original morning sample.}
% END SOURCE BLOCK 137

% BEGIN SOURCE BLOCK 138 paragraph
The mechanism explains destination concentration, but does not uniquely allocate the performance gap to individual decisions. The following recorded accounting locates gains and losses in the original four-draw morning comparison, separately from the eight-block confirmation.
% END SOURCE BLOCK 138

% BEGIN SOURCE BLOCK 139 paragraph
Table~\ref{tab:8} partitions gains and losses over the same 1,744,116 offered tasks per original morning draw: 284,821 gains \ensuremath{-} 12,842 losses = 271,979 net successes. Previously gate-rejected tasks supply most gains, alongside recovered admitted misses; both kinds of losses remain visible. Later queues differ, so these are descriptive transitions, not individual-decision causal effects. The four draws remain the replications \hyperref[source-s13]{[S13]}, \hyperref[source-s14]{[S14]}.
% END SOURCE BLOCK 139

% BEGIN SOURCE BLOCK 140 table
\begingroup
\small
\begin{xltabular}{\textwidth}{@{}YYYYYY@{}}
\caption{Morning ingress\ensuremath{\to}per-task outcome accounting, re-extracted from the authenticated paired results. Gains become successes; losses cease to be successes. All four primary draws are retained; the pilot is excluded.}\label{tab:8} \\
\toprule
Seed & Gain: gate rejection & Gain: admitted miss & Loss: non-admitted & Loss: admitted miss & Net successes \\
\midrule\endfirsthead
\toprule
Seed & Gain: gate rejection & Gain: admitted miss & Loss: non-admitted & Loss: admitted miss & Net successes \\
\midrule\endhead
\bottomrule\endfoot
0 & 53,061 & 18,091 & 395 & 2,945 & 67,812 \\
2 & 64,418 & 19,125 & 496 & 4,144 & 78,903 \\
3 & 52,432 & 18,430 & 105 & 2,808 & 67,949 \\
4 & 41,823 & 17,441 & 48 & 1,901 & 57,315 \\
Total & 211,734 & 73,087 & 1,044 & 11,798 & 271,979 \\
\end{xltabular}
\endgroup
% END SOURCE BLOCK 140

% BEGIN SOURCE BLOCK 141 table
\begingroup
\small
\begin{xltabular}{\textwidth}{@{}YYYYYY@{}}
\caption{Retrospective type-level net additional successes, per-task minus ingress. Each type has the stated offered population in every primary draw; full successes, admissions and failure counts are supplied in Appendix~\ref{app:D}.}\label{tab:9} \\
\toprule
Type / deadline & Offered per draw & Seed 0 & Seed 2 & Seed 3 & Seed 4 \\
\midrule\endfirsthead
\toprule
Type / deadline & Offered per draw & Seed 0 & Seed 2 & Seed 3 & Seed 4 \\
\midrule\endhead
\bottomrule\endfoot
1 / 100 ms & 349,045 & +19,823 & +19,165 & +20,339 & +20,254 \\
2 / 500 ms & 523,060 & +72 & +62 & +6 & +5 \\
3 / 100 ms & 872,011 & +47,917 & +59,676 & +47,604 & +37,056 \\
\end{xltabular}
\endgroup
% END SOURCE BLOCK 141

% BEGIN SOURCE BLOCK 142 paragraph
Type 3 supplies 192,253 net successes (70.7\%), Type 1 79,581 (29.3\%), and Type 2 only 145. Within-type gains span 4.249--6.843, 5.491--5.827 and 0.001--0.014 points, respectively. The two 100 ms types differ substantially in payload and service (\S\,\ref{sec:2.3}). No type loses net successes in any draw, but individual losses persist. In draw 2, Type 1 gate rejections fall by 23,152 while admitted misses rise by 3,987. This supports reduced rejection as an important contributor, without isolating deadlines from communication/service demand or identifying individual-placement causal effects \hyperref[source-s15]{[S15]}.
% END SOURCE BLOCK 142

% BEGIN SOURCE BLOCK 143 heading
\subsection{Bounded state-information and forwarding sensitivities}\label{sec:3.6}
% END SOURCE BLOCK 143

% BEGIN SOURCE BLOCK 144 paragraph
In incident seed 1, fresh per-task attainment was 72.46695\%, versus ingress 72.00328\%. Report ages 100, 500 and 1,000 ms yielded 72.46695\%, 72.46856\% and 72.45510\%; single-draw changes do not establish equivalence or population robustness \hyperref[source-s5]{[S5]}.
% END SOURCE BLOCK 144

% BEGIN SOURCE BLOCK 145 paragraph
Every 100 ms report equalled fresh state over 35,990 non-startup RSU/batch observations: queues had emptied, making that treatment inactive. Reports differed at 500 ms in 79.58\% of observations and at 1,000 ms in all. Live admission and immediate reservations protect this older-report model; stale capacity, delayed acknowledgements or dispersed arrivals would change it. Equal outcomes do not establish harmless communication delay.
% END SOURCE BLOCK 145

% BEGIN SOURCE BLOCK 146 paragraph
Fixed forwarding overheads 0, 1, 2.5, 5 and 10 ms yielded 72.46695\%, 72.44783\%, 72.41882\%, 72.37240\% and 72.28579\% attainment. At 10 ms the advantage remained 0.28251 points, or 36,942 successes over ingress, despite 23,689 extra misses relative to zero overhead \hyperref[source-s6]{[S6]}.
% END SOURCE BLOCK 146

% BEGIN SOURCE BLOCK 147 paragraph
These are saved-record transformations qualified by four short direct probes. Forwarding charges do not feed back into admission, queues or actions in that path; the calculation models neither congestion nor postponed arrivals. A further 104 forwarded tasks had penalty-equal latency ambiguity but were already misses and remained misses under every nonnegative charge. Their latency ambiguity remains unresolved. RQ4 supports the tested information and fixed-cost models, not an extrapolated break-even threshold or network design.
% END SOURCE BLOCK 147

% BEGIN SOURCE BLOCK 148 heading
\subsection{Computational trade-offs and research implementation}\label{sec:3.7}
% END SOURCE BLOCK 148

% BEGIN SOURCE BLOCK 149 table
\begingroup
\small
\begin{xltabular}{\textwidth}{@{}YYYYY@{}}
\caption{Scheduler-only engineering trade-off on Apple M5 CPU, JAX/JAXlib 0.4.30, float32, default runtime threads. Times are milliseconds per five-substep batch; $S$/B denote sparse/busy fixtures. Temporary memory is XLA buffer analysis, not peak RSS.}\label{tab:10} \\
\toprule
Policy & $N=215$ median $S$/B & $N=2{,}488$ median $S$/B & Busy p95, 215/2,488 & Temporary KiB, 215/2,488 \\
\midrule\endfirsthead
\toprule
Policy & $N=215$ median $S$/B & $N=2{,}488$ median $S$/B & Busy p95, 215/2,488 & Temporary KiB, 215/2,488 \\
\midrule\endhead
\bottomrule\endfoot
Ingress & 1.031/1.019 & 6.096/6.113 & 1.103/6.230 & 24.6/294.9 \\
Common-target & 0.952/0.962 & 5.797/5.815 & 1.014/5.935 & 23.8/285.2 \\
Per-task & 0.052/0.055 & 0.501/0.616 & 0.062/0.641 & 11.7/139.1 \\
Round-robin & 0.025/0.021 & 0.223/0.186 & 0.023/0.201 & 11.7/139.1 \\
\end{xltabular}
\endgroup
% END SOURCE BLOCK 149

% BEGIN SOURCE BLOCK 150 paragraph
Per-task was faster than three-pass reconciliation on these CPU fixtures; cyclic selection reduced cost further. Section~\ref{sec:3.4} separately establishes the attainment benefit over cyclic spreading. Sparse padding still traverses $N$ positions. Fixed configuration order and uncontrolled frequency/thermal state limit timing comparisons; full durations and IQRs remain available.
% END SOURCE BLOCK 150

% BEGIN SOURCE BLOCK 151 paragraph
Three vectorised replacements materialise $N\times R$ arrays: $O(KNR)$ work and $O(NR)$ intermediate space with parallel prefix dependence. Per-task also requires $O(KNR)$ work for $R$-way minima, but has a $K\times N$ sequential dependency chain. Round-robin removes minima, giving $O(KN)$ target/predicate work in an indexed-state model. Both retain $R$ queue entries and emit $O(KN)$ diagnostics. Compiler buffer reuse and elimination of unused broadcasts or identical passes affect execution; asymptotic work alone does not predict the measured ordering.
% END SOURCE BLOCK 151

% BEGIN SOURCE BLOCK 152 paragraph
Compilation took 0.058--0.267 seconds per configuration, separately from lowering and execution. Transfer and process peak memory were not measured. The benchmark excludes full-evaluator costs, and host time is never added to historical simulated task latency. No roadside response-time guarantee follows \hyperref[source-s15]{[S15]}.
% END SOURCE BLOCK 152

% BEGIN SOURCE BLOCK 153 paragraph
The research implementation is the trace-driven JAX evaluator, its versioned infrastructure schedulers and the command-line experimental pipeline. Table~\ref{tab:11} shows how implementation choices preserve the meaning of the comparison. The wider TrafficTwin prototype provides interfaces and evidence workflows, but repository membership does not establish that those interfaces executed these campaigns; their commands and outputs are the execution evidence \hyperref[source-s4]{[S4]}, \hyperref[source-s16]{[S16]}.
% END SOURCE BLOCK 153

% BEGIN SOURCE BLOCK 154 table
\begingroup
\small
\begin{xltabular}{\textwidth}{@{}YYY@{}}
\caption{Implementation supporting the research. Verification refers to completed evidence, not additional execution in this editorial revision.}\label{tab:11} \\
\toprule
Inherited starting point & Implemented change and reason & Verification and boundary \\
\midrule\endfirsthead
\toprule
Inherited starting point & Implemented change and reason & Verification and boundary \\
\midrule\endhead
\bottomrule\endfoot
Coarse eligibility / fractional enqueueing & Per-task rejection and full admitted-service enqueueing align tasks and queues & E0 negative count fixture; historical off scoring exception remains \\
Common-target dispatch & Immediate reservations before reselection; reuse the service subkey & Matched service and reconstructed queues detect resampling/double commitment; no physical calibration \\
Existing causal admission & Persistent cyclic targeting changes selection only & Pointer/restart checks and block receipts verify the control population \\
Evaluator outputs / analysis & Separate lifecycle fields and explicit source/seed/completion bindings & Corrupted counts, mismatched inputs and changed hashes are refused; compact checks verify tables only \\
\end{xltabular}
\endgroup
% END SOURCE BLOCK 154

% BEGIN SOURCE BLOCK 155 paragraph
Randy supplied the actor/environment and accounting repair code. Source comments credit Abdulla's accounting questions and clamp reconstruction; the August correspondence records his analysis of the control boundary and interpretation of capacity and placement results. These support specific intellectual contributions beyond coding, while not establishing that they were unaided. The experimental implementations and analyses are verified project outputs. Codex produced substantive drafting and literature work, the retrospective proof/diagnostics, task-type analysis, cyclic comparator, benchmark and confirmation pipeline, and performed its authorised execution/analysis. Personal contributions and checking still require the candidate's confirmation \hyperref[source-s12]{[S12]}, \hyperref[source-s15]{[S15]}, \hyperref[source-s16]{[S16]}.
% END SOURCE BLOCK 155

% BEGIN SOURCE BLOCK 156 paragraph
Two engineering trade-offs limit the result. Aggregate workload/count carry keeps evaluation tractable but does not represent individual departures. Global actual service-work information and immediate reservations make deterministic comparisons inspectable but omit distributed estimation and coordination costs. Separate outcome fields prevent a proposed target being mistaken for execution; they cannot turn either modelling simplification into physical validation. Table~\ref{tab:12} links these limits to the claims they restrict.
% END SOURCE BLOCK 156

% BEGIN SOURCE BLOCK 157 table
\begingroup
\small
\begin{xltabular}{\textwidth}{@{}YYY@{}}
\caption{Validity limits and their consequences for interpretation.}\label{tab:12} \\
\toprule
Limitation & Claim restricted & Consequence / evidence needed \\
\midrule\endfirsthead
\toprule
Limitation & Claim restricted & Consequence / evidence needed \\
\midrule\endhead
\bottomrule\endfoot
Batched arrivals, aggregate departures, backlog-only gate & Physical deadline feasibility & Interpret as evaluator outcomes; validate an event/transfer model against measurements \\
Global actual service work; immediate reservations & Distributed deployability & Measure estimation error and coordination delay before implementation claims \\
Original four fleet draws; new eight joint blocks on one morning trace/actor & Population uncertainty and generalisation & Joint randomness broadens stream coverage; geographical, actor and model uncertainty remain \\
Bundled scenario changes & Morning gain caused by lower density & Separate scenario analyses; factor-controlled evidence needed \\
Legacy off masks differ; E0/E1/E2 originals unavailable & Admission-consistent E1 attainment and gate-only E2b attribution & Historical score intervals remain; measurement impact is not assessable from available records \\
Exact arithmetic versus float32; finite fixtures and unobserved intermediate masks & Numerical prevalence and historical causal decomposition & Two-deadline exact bound proved; constructed rounding discrepancy and reduced adverse example do not establish fleet-frequency effects \\
Single-draw sensitivities & General staleness/forwarding robustness & Describe tested information and fixed-cost models only \\
Unconfirmed backup/access and personal ownership & Unrestricted reproduction and individual achievement & Portable compact checks supplied; raw access and component/AI allocation require confirmation \\
\end{xltabular}
\endgroup
% END SOURCE BLOCK 157

% BEGIN SOURCE BLOCK 158 paragraph
Compact verification regenerates central tables and follows source bindings without the actor or JAX. Raw verification needs separately arranged access; a local copy and checksums are not off-machine preservation. Appendix~\ref{app:A} and the accompanying author notes describe access and attribution boundaries without treating them as additional scientific findings \hyperref[source-s13]{[S13]}, \hyperref[source-s16]{[S16]}.
% END SOURCE BLOCK 158

% BEGIN SOURCE BLOCK 159 heading
\section{Conclusion}\label{sec:4}
% END SOURCE BLOCK 159

% BEGIN SOURCE BLOCK 160 paragraph
TrafficTwin shows that the operational meaning of infrastructure scheduling can determine the direction of a deadline-performance comparison under a frozen vehicle policy. Selecting one common least-workload destination per substep and selecting anew after each admitted task are materially different interventions. Common-target placement underperformed execution at radio ingress, while causal per-task placement exceeded it in the incident study, the original morning replication and the separate joint-randomness confirmation. The contribution is a controlled and explainable implementation study of established scheduling ideas.
% END SOURCE BLOCK 160

% BEGIN SOURCE BLOCK 161 paragraph
RQ1 supplied the measurement foundation and a caution about capacity. Counting all offered tasks keeps rejection visible, while separate admission and outcome fields expose inconsistencies that aggregate enqueueing can conceal. Increasing the waiting-room limit did not establish an improvement in the archived E1 score: its interval is statistically inconclusive. Interpretation as admission-consistent deadline attainment also retains the unquantified historical scoring-mask exception. Reproducing that interval from summaries resolves neither the missing task-level check nor the possibility of feedback through energy and later observations. A larger waiting room is therefore neither evidence of faster service nor a demonstrated solution to deadline failure.
% END SOURCE BLOCK 161

% BEGIN SOURCE BLOCK 162 paragraph
For RQ2, the choice of comparator was decisive. The exploratory admission decomposition showed why strongest-link execution with the gate was the relevant reference, while its eligibility/scoring differences limit a gate-only interpretation. Subsequent matched incident comparisons reversed the ranking when the dispatch implementation changed. The adaptive per-task extension reused examined controls; the later morning studies provide separate evidence rather than retroactively making that extension confirmatory.
% END SOURCE BLOCK 162

% BEGIN SOURCE BLOCK 163 paragraph
RQ3 is answered by two distinct morning samples. The original four fleet draws, excluding the inspected pilot, reproduced the reversal under their simultaneous interval family. The new eight joint fleet/evaluator blocks reproduced it again: per-task minus ingress was +4.137 percentage points and common-target minus ingress \ensuremath{-}3.504. Per-task also exceeded causal cyclic spreading by +0.631 points, with its simultaneous interval entirely positive. Workload information therefore added benefit beyond spreading in this study. All actions matched, despite small observation/logit differences in one block; equality was observed rather than forced. The practical gain over round-robin accompanies higher scheduler-only cost on separate CPU fixtures, not measured roadside or full-system latency.
% END SOURCE BLOCK 163

% BEGIN SOURCE BLOCK 164 paragraph
Retrospective explanation sharpens, but does not uniquely decompose, those findings. The exact-model result rules out unrestricted fixed-target reconciliation failure under the studied two-deadline assumptions; gate, service and drain semantics explain recurring common destinations. Float32 boundaries and the reduced adverse example prevent treating that result as universal executable equivalence or least-workload optimality. Original four-draw task accounting shows that reduced rejection supplies most gains, with additional recovery of admitted misses and smaller losses elsewhere. These are matched outcome transitions, not individual-decision causal effects.
% END SOURCE BLOCK 164

% BEGIN SOURCE BLOCK 165 paragraph
RQ4 exposes the importance of checking whether a treatment changes information: the 100 ms reports equalled fresh state. Live admission and reservations also protect the tested older-report model. The single-draw forwarding sensitivity retained an advantage through 10 ms, but fixed overhead does not model network congestion.
% END SOURCE BLOCK 165

% BEGIN SOURCE BLOCK 166 paragraph
The research process consequently shifted from comparing labels to checking denominators, admission, reservations and controls. Future work should prioritise event-level service and communication measurements to test physical timing, imperfect workload information and acknowledgement to test coordination assumptions, and different traces or actors to test transfer beyond this setting. Each would resolve a distinct uncertainty rather than simply enlarge the experiment count. The established result remains bounded: within the studied evaluator and matched populations, implementation semantics reverse the ingress comparison, and workload-aware causal placement adds a measured attainment benefit over the tested spreading rule.
% END SOURCE BLOCK 166

% BEGIN SOURCE BLOCK 167 bibliography_open
\begin{thebibliography}{99}
% END SOURCE BLOCK 167

% BEGIN SOURCE BLOCK 168 bibliography_item
\bibitem{ref1} Y. Mao, C. You, J. Zhang, K. Huang and K. B. Letaief. ``A Survey on Mobile Edge Computing: The Communication Perspective.'' \emph{IEEE Communications Surveys \& Tutorials}, 19(4), 2322--2358, 2017. DOI: 10.1109/COMST.2017.2745201. \href{\detokenize{https://arxiv.org/abs/1701.01090}}{Author manuscript}.
% END SOURCE BLOCK 168

% BEGIN SOURCE BLOCK 169 bibliography_item
\bibitem{ref2} M. Harchol-Balter, M. E. Crovella and C. D. Murta. ``On Choosing a Task Assignment Policy for a Distributed Server System.'' \emph{Computer Performance Evaluation (TOOLS 1998)}, LNCS 1469, 231--242, 1998. DOI: 10.1007/3-540-68061-6\_19. \href{\detokenize{https://www.cs.cmu.edu/~harchol/Papers/tools.pdf}}{Author manuscript}.
% END SOURCE BLOCK 169

% BEGIN SOURCE BLOCK 170 bibliography_item
\bibitem{ref3} J. Schulman, F. Wolski, P. Dhariwal, A. Radford and O. Klimov. ``Proximal Policy Optimization Algorithms.'' arXiv:1707.06347, 2017. \href{\detokenize{https://arxiv.org/abs/1707.06347}}{Author manuscript}.
% END SOURCE BLOCK 170

% BEGIN SOURCE BLOCK 171 bibliography_item
\bibitem{ref4} C. Yu, A. Velu, E. Vinitsky, J. Gao, Y. Wang, A. Bayen and Y. Wu. ``The Surprising Effectiveness of PPO in Cooperative Multi-Agent Games.'' \emph{Advances in Neural Information Processing Systems}, 35, 24611--24624, 2022. \href{\detokenize{https://papers.nips.cc/paper/2022/file/9c1535a02f0ce079433344e14d910597-Paper-Datasets_and_Benchmarks.pdf}}{Proceedings paper}.
% END SOURCE BLOCK 171

% BEGIN SOURCE BLOCK 172 bibliography_item
\bibitem{ref5} P. Henderson, R. Islam, P. Bachman, J. Pineau, D. Precup and D. Meger. ``Deep Reinforcement Learning That Matters.'' \emph{Proceedings of the AAAI Conference on Artificial Intelligence}, 32(1), 2018. DOI: 10.1609/aaai.v32i1.11694. \href{\detokenize{https://ojs.aaai.org/index.php/AAAI/article/view/11694}}{Publisher record}.
% END SOURCE BLOCK 172

% BEGIN SOURCE BLOCK 173 bibliography_item
\bibitem{ref6} R. Agarwal, M. Schwarzer, P. S. Castro, A. Courville and M. G. Bellemare. ``Deep Reinforcement Learning at the Edge of the Statistical Precipice.'' \emph{Advances in Neural Information Processing Systems}, 34, 29304--29320, 2021. \href{\detokenize{https://proceedings.neurips.cc/paper/2021/hash/f514cec81cb148559cf475e7426eed5e-Abstract.html}}{Proceedings record}.
% END SOURCE BLOCK 173

% BEGIN SOURCE BLOCK 174 bibliography_item
\bibitem{ref7} R. Gorsane, O. Mahjoub, R. de Kock, R. Dubb, S. Singh and A. Pretorius. ``Towards a Standardised Performance Evaluation Protocol for Cooperative MARL.'' \emph{Advances in Neural Information Processing Systems}, 35, 5510--5521, 2022. \href{\detokenize{https://arxiv.org/abs/2209.10485}}{Author manuscript}.
% END SOURCE BLOCK 174

% BEGIN SOURCE BLOCK 175 bibliography_item
\bibitem{ref8} P. Alvarez Lopez, M. Behrisch, L. Bieker-Walz, J. Erdmann, Y.-P. Flötteröd, R. Hilbrich, L. Lücken, J. Rummel, P. Wagner and E. Wießner. ``Microscopic Traffic Simulation using SUMO.'' \emph{21st International Conference on Intelligent Transportation Systems (ITSC)}, 2575--2582, 2018. DOI: 10.1109/ITSC.2018.8569938. \href{\detokenize{https://elib.dlr.de/127994/}}{DLR author repository}.
% END SOURCE BLOCK 175

% BEGIN SOURCE BLOCK 176 bibliography_item
\bibitem{ref9} K. Ousterhout, P. Wendell, M. Zaharia and I. Stoica. ``Sparrow: Distributed, Low Latency Scheduling.'' \emph{24th ACM Symposium on Operating Systems Principles}, 69--84, 2013. DOI: 10.1145/2517349.2522716. \href{\detokenize{https://sigops.org/s/conferences/sosp/2013/papers/p69-ousterhout.pdf}}{Proceedings paper}.
% END SOURCE BLOCK 176

% BEGIN SOURCE BLOCK 177 bibliography_item
\bibitem{ref10} S. Vargaftik, I. Keslassy and A. Orda. ``LSQ: Load Balancing in Large-Scale Heterogeneous Systems With Multiple Dispatchers.'' \emph{IEEE/ACM Transactions on Networking}, 28(3), 1186--1198, 2020. DOI: 10.1109/TNET.2020.2980061. \href{\detokenize{https://webee.technion.ac.il/~isaac/p/ton20_lsq.pdf}}{Author manuscript}.
% END SOURCE BLOCK 177

% BEGIN SOURCE BLOCK 178 bibliography_item
\bibitem{ref11} J. Niño-Mora. ``Admission and routing of soft real-time jobs to multiclusters: Design and comparison of index policies.'' \emph{Computers \& Operations Research}, 39, 3431--3444, 2012. DOI: 10.1016/j.cor.2012.05.004. \href{\detokenize{https://arxiv.org/abs/2207.12815}}{Author manuscript, deposited 2022}.
% END SOURCE BLOCK 178

% BEGIN SOURCE BLOCK 179 bibliography_item
\bibitem{ref12} R. G. Sargent. ``Verification and Validation of Simulation Models.'' \emph{Proceedings of the 2011 Winter Simulation Conference}, 183--198, 2011. \href{\detokenize{https://www.informs-sim.org/wsc11papers/016.pdf}}{Proceedings paper}.
% END SOURCE BLOCK 179

% BEGIN SOURCE BLOCK 180 bibliography_item
\bibitem{ref13} L. Ying, R. Srikant and X. Kang. ``The Power of Slightly More than One Sample in Randomized Load Balancing.'' \emph{IEEE INFOCOM}, 1131--1139, 2015. \href{\detokenize{https://doi.org/10.1109/INFOCOM.2015.7218487}}{doi:10.1109/INFOCOM.2015.7218487}. \href{\detokenize{https://bpb-us-w2.wpmucdn.com/sites.coecis.cornell.edu/dist/9/287/files/2019/08/Srikant-1-YinSriKan14tech.pdf}}{Author technical manuscript}.
% END SOURCE BLOCK 180

% BEGIN SOURCE BLOCK 181 bibliography_item
\bibitem{ref14} JAX authors. ``Benchmarking JAX code.'' Official JAX documentation, accessed 8 September 2026. \href{\detokenize{https://docs.jax.dev/en/latest/benchmarking.html}}{Benchmarking guidance}. Used for timing procedure, not scheduler efficacy.
% END SOURCE BLOCK 181

% BEGIN SOURCE BLOCK 182 bibliography_close
\end{thebibliography}
% END SOURCE BLOCK 182

% BEGIN SOURCE BLOCK 183 appendix_open
\appendix
\renewcommand{\thetable}{\Alph{section}\arabic{table}}
% END SOURCE BLOCK 183

% BEGIN SOURCE BLOCK 184 heading
\appendixsection{app:A}{Evidence sources and reproducibility}
% END SOURCE BLOCK 184

% BEGIN SOURCE BLOCK 185 paragraph
The sources below are project evidence, separate from scholarly references. Relative links resolve in the repository checkout. The \href{\detokenize{https://github.com/Abdulla4akash/traffictwin/tree/04f3b6a95c7bc06c80ed95b54762f12861bba183}}{historical baseline} and later packages retain their respective source identities. The current \href{\detokenize{CLAIM_SOURCE_MAP.md}}{claim-to-source map} records support and limits.
% END SOURCE BLOCK 185

% BEGIN SOURCE BLOCK 186 anchor
\phantomsection\label{source-s0}
% END SOURCE BLOCK 186

% BEGIN SOURCE BLOCK 187 paragraph
\textbf{S0 --- Accounting validity.} \href{\detokenize{../../evaluation/vec_followup_2026-09-07/historical_e0_e2d/experiments/E0/README.md}}{E0 report and evidence}. Full evidence commit \texttt{29d8945862b7df1bd5d7879ba1d980a388f6be0d}.
% END SOURCE BLOCK 187

% BEGIN SOURCE BLOCK 188 anchor
\phantomsection\label{source-s1}
% END SOURCE BLOCK 188

% BEGIN SOURCE BLOCK 189 paragraph
\textbf{S1 --- Waiting-room capacity.} \href{\detokenize{../../evaluation/vec_followup_2026-09-07/historical_e0_e2d/experiments/E1/README.md}}{E1 report}, with five-draw comparison and validation records.
% END SOURCE BLOCK 189

% BEGIN SOURCE BLOCK 190 anchor
\phantomsection\label{source-s2}
% END SOURCE BLOCK 190

% BEGIN SOURCE BLOCK 191 paragraph
\textbf{S2 --- Exploratory admission/placement decomposition.} \href{\detokenize{../../evaluation/vec_followup_2026-09-07/historical_e0_e2d/experiments/E2b/README.md}}{E2b report and factorial evidence}.
% END SOURCE BLOCK 191

% BEGIN SOURCE BLOCK 192 anchor
\phantomsection\label{source-s3}
% END SOURCE BLOCK 192

% BEGIN SOURCE BLOCK 193 paragraph
\textbf{S3 --- Incident replication and reversal.} \href{\detokenize{../../evaluation/vec_followup_2026-09-07/historical_e0_e2d/experiments/E2c/README.md}}{E2c report} and \href{\detokenize{../../evaluation/vec_followup_2026-09-07/historical_e0_e2d/experiments/E2d/README.md}}{E2d report}. Frozen E2d evaluator commit \texttt{2f63706f46319433a2ba3af1df97afd0e56a95d1}; exact E2c controls reused.
% END SOURCE BLOCK 193

% BEGIN SOURCE BLOCK 194 anchor
\phantomsection\label{source-s4}
% END SOURCE BLOCK 194

% BEGIN SOURCE BLOCK 195 paragraph
\textbf{S4 --- Experimental architecture and implementation.} \href{\detokenize{../../evaluation/vec_followup_2026-09-07/historical_e0_e2d/docs/METHODOLOGY.md}}{Frozen method}, \href{\detokenize{../../evaluation/vec_followup_2026-09-07/frozen_evaluator/eval/e2d_per_task_placement.py}}{per-task source}, \href{\detokenize{../../evaluation/vec_followup_2026-09-07/frozen_evaluator/eval/eval_sumo_stage1_mc.py}}{evaluator source}, \href{\detokenize{../empirical_extension_2026-09-08/experimental/vendor/vec_jax.py}}{environment source}, and \href{\detokenize{../../correspondence/sandra_randy_vec_progress_email_thread_2026-08-18.md}}{implementation correspondence}, PDF pages 1--2 for Randy's confirmations.
% END SOURCE BLOCK 195

% BEGIN SOURCE BLOCK 196 anchor
\phantomsection\label{source-s5}
% END SOURCE BLOCK 196

% BEGIN SOURCE BLOCK 197 paragraph
\textbf{S5 --- State-information sensitivity.} \href{\detokenize{../../evaluation/vec_followup_2026-09-07/frozen_evaluator/docs/RSU_STATE_DELAY.md}}{Timing contract}, \href{\detokenize{../../evaluation/vec_followup_2026-09-07/state-delay-pilot-2026-09-07/comparison.csv}}{five-cell comparison} and \href{\detokenize{../../evaluation/vec_followup_2026-09-07/state-delay-pilot-2026-09-07/mechanism_audit.json}}{report-state audit}.
% END SOURCE BLOCK 197

% BEGIN SOURCE BLOCK 198 anchor
\phantomsection\label{source-s6}
% END SOURCE BLOCK 198

% BEGIN SOURCE BLOCK 199 paragraph
\textbf{S6 --- Forwarding sensitivity.} \href{\detokenize{../../evaluation/vec_followup_2026-09-07/forwarding-sensitivity-2026-09-07/qualification.json}}{Qualification}, \href{\detokenize{../../evaluation/vec_followup_2026-09-07/forwarding-sensitivity-2026-09-07/forwarding_sensitivity.csv}}{outcomes} and \href{\detokenize{../../evaluation/vec_followup_2026-09-07/forwarding-sensitivity-2026-09-07/analysis_validation.json}}{analysis validation}.
% END SOURCE BLOCK 199

% BEGIN SOURCE BLOCK 200 anchor
\phantomsection\label{source-s7}
% END SOURCE BLOCK 200

% BEGIN SOURCE BLOCK 201 paragraph
\textbf{S7 --- Morning exploratory pilot.} \href{\detokenize{../../evaluation/vec_followup_2026-09-07/generalisation-pilot-2026-09-07/RESULTS.md}}{Results} and \href{\detokenize{../../evaluation/vec_followup_2026-09-07/generalisation-pilot-2026-09-07/input_validation.json}}{canonical input validation}.
% END SOURCE BLOCK 201

% BEGIN SOURCE BLOCK 202 anchor
\phantomsection\label{source-s8}
% END SOURCE BLOCK 202

% BEGIN SOURCE BLOCK 203 paragraph
\textbf{S8 --- Morning primary replication.} \href{\detokenize{../../evaluation/vec_followup_2026-09-07/generalisation-replication-2026-09-07/PROTOCOL.md}}{Protocol}, \href{\detokenize{../../evaluation/vec_followup_2026-09-07/generalisation-replication-2026-09-07/manifest.json}}{manifest}, \href{\detokenize{../../evaluation/vec_followup_2026-09-07/generalisation-replication-2026-09-07/paired_differences.csv}}{paired differences} and \href{\detokenize{../../evaluation/vec_followup_2026-09-07/generalisation-replication-2026-09-07/paired_intervals.csv}}{paired intervals}. Full follow-up evaluator commit \texttt{908bd10f86542de94fc38af90dd56c2ccc08cf9b}.
% END SOURCE BLOCK 203

% BEGIN SOURCE BLOCK 204 anchor
\phantomsection\label{source-s9}
% END SOURCE BLOCK 204

% BEGIN SOURCE BLOCK 205 paragraph
\textbf{S9 --- Five-RSU mechanism audit.} \href{\detokenize{../../evaluation/common_target_mechanism_audit_2026-09-07/evidence/REPORT.md}}{Completed post-hoc report} and \href{\detokenize{../../evaluation/common_target_mechanism_audit_2026-09-07/evidence/audit_results.json}}{audit results with input identities}.
% END SOURCE BLOCK 205

% BEGIN SOURCE BLOCK 206 anchor
\phantomsection\label{source-s10}
% END SOURCE BLOCK 206

% BEGIN SOURCE BLOCK 207 paragraph
\textbf{S10 --- Unrun falsification proposal.} \href{\detokenize{../../evaluation/common_target_mechanism_audit_2026-09-07/TIEBREAK_PREFIX_PROPOSAL.md}}{Rotating-tie-break prefix protocol}.
% END SOURCE BLOCK 207

% BEGIN SOURCE BLOCK 208 anchor
\phantomsection\label{source-s11}
% END SOURCE BLOCK 208

% BEGIN SOURCE BLOCK 209 paragraph
\textbf{S11 --- Artifact and preservation.} \href{\detokenize{../../architecture_current_release.md}}{Product architecture}, \href{\detokenize{../../quality/final_release_status_20260815.md}}{release/execution boundary}, \href{\detokenize{../../evaluation/vec_followup_2026-09-07/ARCHIVE_INVENTORY.json}}{evidence archive inventory} and \href{\detokenize{../../evaluation/vec_followup_2026-09-07/COMPACT_VERIFICATION.json}}{compact verification receipt}. The separate E3 Dynamic Resource V2 campaign remains unexecuted.
% END SOURCE BLOCK 209

% BEGIN SOURCE BLOCK 210 anchor
\phantomsection\label{source-s12}
% END SOURCE BLOCK 210

% BEGIN SOURCE BLOCK 211 paragraph
\textbf{S12 --- Contribution and repair history.} Read-only inspection of \texttt{vec\_env} commits \href{\detokenize{https://github.com/Abdulla4akash/vec_env/commit/4cb7c06aab1b455d68f71ea9143172463784dea9}}{4cb7c06}, \href{\detokenize{https://github.com/Abdulla4akash/vec_env/commit/0f01f4d2082d3e8b735e74a873095ab8eeba37cc}}{0f01f4d} and \href{\detokenize{https://github.com/Abdulla4akash/vec_env/commit/2f63706f46319433a2ba3af1df97afd0e56a95d1}}{2f63706}, and TrafficTwin \href{\detokenize{https://github.com/Abdulla4akash/traffictwin/blob/29d8945862b7df1bd5d7879ba1d980a388f6be0d/tests/test_validate_e0_smoke.py}}{E0 validation tests at 29d894}. Source comments explicitly credit Abdulla's questions; commit authorship does not establish independent human authorship.
% END SOURCE BLOCK 211

% BEGIN SOURCE BLOCK 212 anchor
\phantomsection\label{source-s13}
% END SOURCE BLOCK 212

% BEGIN SOURCE BLOCK 213 paragraph
\textbf{S13 --- Retrospective gap closure, 8 September 2026.} The \href{\detokenize{../gap_closure_2026-09-08/verification/README.md}}{verification entry point} documents the field contract, source bindings, 87-file authentication, 19-run audit, \href{\detokenize{../gap_closure_2026-09-08/verification/results/OUTCOME_TRANSITIONS.json}}{23 paired transitions}, \href{\detokenize{../gap_closure_2026-09-08/verification/results/KERNEL_RESULTS.json}}{440-case reference} and production comparison. These completed checks and their historical sources remain unchanged.
% END SOURCE BLOCK 213

% BEGIN SOURCE BLOCK 214 anchor
\phantomsection\label{source-s14}
% END SOURCE BLOCK 214

% BEGIN SOURCE BLOCK 215 paragraph
\textbf{S14 --- Production-domain mechanism analysis.} \href{\detokenize{../production_mechanism_2026-09-08/analysis/MATHEMATICS.md}}{Model and proof}, \href{\detokenize{../production_mechanism_2026-09-08/analysis/PROTOCOL.md}}{initial diagnostic protocol}, \href{\detokenize{../production_mechanism_2026-09-08/analysis/FOLLOWUP_PROTOCOL.md}}{bounded follow-up protocol}, \href{\detokenize{../production_mechanism_2026-09-08/analysis/diagnostics.py}}{executable diagnostics}, \href{\detokenize{../production_mechanism_2026-09-08/analysis/results/DIAGNOSTICS.json}}{full results}, \href{\detokenize{../production_mechanism_2026-09-08/analysis/results/NUMERICAL_FOLLOWUP.json}}{numerical follow-up}, \href{\detokenize{../production_mechanism_2026-09-08/analysis/results/ADVERSE_PRODUCTION.json}}{adverse example and helper checks}, \href{\detokenize{../production_mechanism_2026-09-08/analysis/results/DIMENSION_CHECK.json}}{dimension checks} and \href{\detokenize{../production_mechanism_2026-09-08/analysis/results/OUTCOME_ACCOUNTING.json}}{outcome accounting}. This is retrospective analysis added after the completed experiments and gap closure.
% END SOURCE BLOCK 215

% BEGIN SOURCE BLOCK 216 anchor
\phantomsection\label{source-s15}
% END SOURCE BLOCK 216

% BEGIN SOURCE BLOCK 217 paragraph
\textbf{S15 --- Retrospective empirical/engineering extension.} \href{\detokenize{../empirical_extension_2026-09-08/analysis/PROTOCOL.md}}{Declared analysis protocol}, \href{\detokenize{../empirical_extension_2026-09-08/analysis/results/TASK_TYPES.json}}{task-type counts and bindings}, \href{\detokenize{../empirical_extension_2026-09-08/analysis/benchmark.py}}{benchmark implementation} and \href{\detokenize{../empirical_extension_2026-09-08/analysis/results/BENCHMARK.json}}{measured results}, \href{\detokenize{../empirical_extension_2026-09-08/experimental/round_robin.py}}{cyclic implementation}, \href{\detokenize{../empirical_extension_2026-09-08/analysis/results/COMPATIBILITY_TESTS.json}}{compatibility evidence}, \href{\detokenize{../empirical_extension_2026-09-08/confirmation/PROTOCOL.md}}{sealed confirmation protocol}, \href{\detokenize{../empirical_extension_2026-09-08/confirmation/DRY_RUN.json}}{dry-run receipt}, and \href{\detokenize{../empirical_extension_2026-09-08/README.md}}{verification route}. That original preparation package remains disabled and unchanged; S16 records the separately authorised execution.
% END SOURCE BLOCK 217

% BEGIN SOURCE BLOCK 218 anchor
\phantomsection\label{source-s16}
% END SOURCE BLOCK 218

% BEGIN SOURCE BLOCK 219 paragraph
\textbf{S16 --- Authorised joint-randomness confirmation.} \href{\detokenize{../joint_confirmation_2026-09-08/confirmation/EXECUTION_AMENDMENT.md}}{Execution amendment}, \href{\detokenize{../joint_confirmation_2026-09-08/confirmation/PROTOCOL_EXECUTION.md}}{qualified execution protocol}, \href{\detokenize{../joint_confirmation_2026-09-08/confirmation/SEALED_EXECUTION.json}}{execution seal}, \href{\detokenize{../joint_confirmation_2026-09-08/experimental/evaluator_v2.py}}{versioned evaluator}, \href{\detokenize{../joint_confirmation_2026-09-08/confirmation/runner.py}}{runner and completion gate}, \href{\detokenize{../joint_confirmation_2026-09-08/confirmation/validation.py}}{task/queue/block validation}, \href{\detokenize{../joint_confirmation_2026-09-08/confirmation/analyse.py}}{sealed analysis}, \href{\detokenize{../joint_confirmation_2026-09-08/evidence/QUALIFICATION.json}}{qualification receipt}, \href{\detokenize{../joint_confirmation_2026-09-08/evidence/QUALIFICATION_FOLLOWUP.json}}{prelaunch safeguard checks}, and \href{\detokenize{README.md}}{verification/access route}. \href{\detokenize{../joint_confirmation_2026-09-08/evidence/ANALYSIS.json}}{Completed primary analysis}, \href{\detokenize{../joint_confirmation_2026-09-08/evidence/CELL_RESULTS.csv}}{all cell counts}, \href{\detokenize{../joint_confirmation_2026-09-08/evidence/BLOCK_CONTROLS.json}}{block controls} and \href{\detokenize{../joint_confirmation_2026-09-08/evidence/RAW_INVENTORY.json}}{raw inventory} bind the new results. This study is separate from the original four-draw sample, pilot and retrospective analyses.
% END SOURCE BLOCK 219

% BEGIN SOURCE BLOCK 220 paragraph
The runtime is CPython 3.11.15, JAX/JAXlib 0.4.30 and NumPy 1.26.4, CPU, x64 disabled. Historical manifests retain their own identities. Full reproduction also needs the frozen actor, traces and authorised inputs. E0/E1/E2 originals remain unavailable following author-reported deletion, with no known backup; other historical raw availability is unestablished.
% END SOURCE BLOCK 220

% BEGIN SOURCE BLOCK 221 paragraph
\textbf{Contribution record.} The accompanying \href{\detokenize{AUTHOR_INPUTS.md}}{author notes} identify the personal decisions, checking and assessment-specific disclosure still to confirm. Recorded supplied and AI-assisted work is stated in \S\,\ref{sec:3.7} and Appendix~\ref{app:C}.
% END SOURCE BLOCK 221

% BEGIN SOURCE BLOCK 222 heading
\appendixsection{app:B}{Audit coverage and diagnostic boundaries}
% END SOURCE BLOCK 222

% BEGIN SOURCE BLOCK 223 paragraph
Table~\ref{tab:B1} retains the completed gap-closure audit coverage. Those raw checks and joins were not repeated for this revision.
% END SOURCE BLOCK 223

% BEGIN SOURCE BLOCK 224 table
\begingroup
\small
\begin{xltabular}{\textwidth}{@{}YYYY@{}}
\caption{Historical availability and retained gap-closure audit coverage, 8 September 2026. Historical receipts remain unchanged.}\label{tab:B1} \\
\toprule
Records & Available evidence / coverage & Completed gap-closure checks & Consequence \\
\midrule\endfirsthead
\toprule
Records & Available evidence / coverage & Completed gap-closure checks & Consequence \\
\midrule\endhead
\bottomrule\endfoot
E0: two ten-step repeats and full seed 0 & Surviving receipts, manifests, sources; full receipt records 834,120 cap rejections & Not assessable from available records & Original raw outputs unavailable; deletion reported by author; no known backup \\
E1: five seeds \ensuremath{\times} three caps & Fifteen run summaries/receipts; seed-0 middle cap reuses E0 & Not assessable from available records & Original interval reproduced from summaries; scoring impact unresolved \\
E2 off/jsq/dla seed 0 and E2b reused cells & Hash-bound summaries and receipts; E2b reuses E2 exactly & Not assessable from available records & No new cell or historical task join inferred from reuse \\
E2b new ingress / E2c / E2d originals & Compact sources/records survive; original raw availability not established & Not assessable from available records & No deletion or fresh raw-validation claim extended to these studies \\
September morning primary & Twelve full runs, seeds 0/2/3/4; 1,744,116 offered per cell & Zero discrepancies in each inspected run & Supports these gate-enabled records; excludes pilot from inference \\
September morning pilot & Two full seed-1 runs, 1,744,116 offered each & Zero discrepancies in each inspected run & Remains exploratory, outside primary inference \\
September incident sensitivity pilot & Five full seed-1 runs, 13,076,234 offered each & Zero discrepancies in each inspected run & Does not replace historical E2d records \\
Other September arrays & Remaining inventory entries authenticated only & Integrity-only in the prior audit & Qualification/preflight files are not additional fleet replications \\
\end{xltabular}
\endgroup
% END SOURCE BLOCK 224

% BEGIN SOURCE BLOCK 225 paragraph
The per-run JSON records preserve offered/admitted counts, explicit rejection and unavailability categories, false-success and latency discrepancy counts, penalty-equal ambiguity, summary reconciliation and source bindings. A rejected task's finite latency is not necessarily a false success. For the 19 available full runs, both false-success and non-penalty rejection-latency counts are zero, and fixed-history admission-consistent scores equal their archived values. Historical unavailable cells are labelled ``not assessable,'' never zero. Receipt checks are reported as past observations, not freshly rerun validation.
% END SOURCE BLOCK 225

% BEGIN SOURCE BLOCK 226 paragraph
Table~\ref{tab:B2} details primary morning ingress\ensuremath{\to}per-task task accounting. ``Both'' columns partition the same offered population; no task-level confidence intervals are constructed.
% END SOURCE BLOCK 226

% BEGIN SOURCE BLOCK 227 table
\begingroup
\small
\begin{xltabular}{\textwidth}{@{}YYYYYY@{}}
\caption{Descriptive task transitions from ingress to per-task placement; four primary morning draws, pilot excluded.}\label{tab:B2} \\
\toprule
Fleet seed & Failure \ensuremath{\to} success & Success \ensuremath{\to} failure & Success both & Failure both & Net successes \\
\midrule\endfirsthead
\toprule
Fleet seed & Failure \ensuremath{\to} success & Success \ensuremath{\to} failure & Success both & Failure both & Net successes \\
\midrule\endhead
\bottomrule\endfoot
0 & 71,152 & 3,340 & 1,545,724 & 123,900 & 67,812 \\
2 & 83,543 & 4,640 & 1,522,085 & 133,848 & 78,903 \\
3 & 70,862 & 2,913 & 1,550,789 & 119,552 & 67,949 \\
4 & 59,264 & 1,949 & 1,569,719 & 113,184 & 57,315 \\
\end{xltabular}
\endgroup
% END SOURCE BLOCK 227

% BEGIN SOURCE BLOCK 228 paragraph
All 23 eligible within-study pairs are retained in the machine-readable record. Task identity uses matched trace-row/substep/slot coordinates and verified stream identities; sampled sizes are source-derived correspondence, not independently observed fields. Later queue divergence prevents these descriptive transitions from identifying a unique causal contribution of individual placements.
% END SOURCE BLOCK 228

% BEGIN SOURCE BLOCK 229 paragraph
The scalar suite includes 440 primary cases and ten deletion trials, below its 500-case limit. The five-task case is minimal under single-candidate deletion, not a globally minimal counterexample over every numeric parameter. A's retained mask remains unchanged by the diagnostic code: the fourth replacement is inspected only. C's repeated complete passes restart from the same state. That scalar exercise performed no full-evaluator prefix, actor inference, SUMO run, campaign or retraining; the new authorised campaign is separate.
% END SOURCE BLOCK 229

% BEGIN SOURCE BLOCK 230 heading
\appendixsection{app:C}{Exact-model result and numerical boundary}
% END SOURCE BLOCK 230

% BEGIN SOURCE BLOCK 231 paragraph
The following full model-specific proof is reproduced from S14, preserving its assumptions and numerical boundary. Codex produced this retrospective analysis; independent student authorship or completed personal verification is not asserted. The existing diagnostic results follow the proof; their suite has not been repeated.
% END SOURCE BLOCK 231

% BEGIN SOURCE BLOCK 232 paragraph
This is a retrospective deduction about the stated TrafficTwin admission model. It is not a new scheduling family, a full-system experiment or a claim of independent student authorship. The frozen code is not altered. Numerical findings and executable inputs accompany this note.
% END SOURCE BLOCK 232

% BEGIN SOURCE BLOCK 233 heading
\subsection{Model and objects that must remain separate}\label{sec:C.1}
% END SOURCE BLOCK 233

% BEGIN SOURCE BLOCK 234 paragraph
Candidates have a fixed order $i=1,\ldots,n$ and fixed destination $r_i$. At each destination $r$, initial workload $W_r\geq0$, integer task count $Q_r\geq0$ and capacity $C_r$ are fixed throughout reconciliation. $E_i$ is fixed active/radio/coarse-capacity eligibility. Candidate service $s_i\geq0$ and deadline $D_i$ are finite. There is no intervening service drain. All sums here are exact.
% END SOURCE BLOCK 234

% BEGIN SOURCE BLOCK 235 paragraph
For mask $M\leq E$, define $O_i(M)=\sum_{\substack{j<i\\r_j=r_i}}s_jM_j$ and $H_i(M)=\sum_{\substack{j<i\\r_j=r_i}}M_j$. The map is
% END SOURCE BLOCK 235

% BEGIN SOURCE BLOCK 236 equation
\begin{equation}\tag{C1}\label{eq:C1}
 F_i(M)=E_i\,\mathbf{1}\{W_{r_i}+O_i(M)<D_i\}\,\mathbf{1}\{Q_{r_i}+H_i(M)<C_{r_i}\}.
\end{equation}
% END SOURCE BLOCK 236

% BEGIN SOURCE BLOCK 237 paragraph
Start $M^{0}=E$; the historical algorithm retains $M^{3}=F^{3}(E)$, then recomputes $O(M^{3})$. Its gate-failure label is $E_i\mathbf{1}\{W+O_i(M^{3})\geq D_i\}$; its capacity label covers other eligible non-admissions, even when the retained mask is not fixed. Enqueueing follows $M^{3}$. Diagnostic success is $M_i^{3}\mathbf{1}\{W+O_i(M^{3})+s_i\leq D_i\}$, with zero transfer latency for these fixtures. Production latency additionally contains the source-defined communication/forwarding terms. Neither admission nor a rejection label alone determines the final deadline outcome.
% END SOURCE BLOCK 237

% BEGIN SOURCE BLOCK 238 paragraph
The causal fixed-target scan computes $b_i=F_i(b)$ in order, reserving $s_i$ and one count immediately when admitted. It uses the same $E$, destinations and input arrays. This scan is B in a one-substep comparison. Across multiple substeps B replays A's entire target sequence; it is a target-conditioned diagnostic, not automatically an autonomous deployment policy.
% END SOURCE BLOCK 238

% BEGIN SOURCE BLOCK 239 heading
\subsection{Existence, uniqueness and candidate-count bound}\label{sec:C.2}
% END SOURCE BLOCK 239

% BEGIN SOURCE BLOCK 240 proof_open
\begin{proof}
% END SOURCE BLOCK 240

% BEGIN SOURCE BLOCK 241 paragraph
$F_i$ depends only on earlier candidates at the same destination. Define $b_1$ from its constant predicate, then $b_2$ from $b_1$, and so on. This constructs a fixed point. If two fixed points differed, their earliest differing candidate would see identical preceding inputs and hence could not differ. Thus $b$ is unique and is exactly the causal scan.
% END SOURCE BLOCK 241

% BEGIN SOURCE BLOCK 242 paragraph
From any initial mask, after $t$ simultaneous replacements the first $t$ eligible candidates at each fixed destination have their final values: induction uses the already-correct preceding values. Therefore at most $n_r$ replacements suffice at destination $r$, and at most $\max_r n_r$ globally. Ineligible candidates can be removed because their mask is fixed at zero. Multiple fixed destinations form disjoint dependency chains. This argument alone does not justify three replacements for an arbitrary candidate population.
% END SOURCE BLOCK 242

% BEGIN SOURCE BLOCK 243 proof_close
\end{proof}
% END SOURCE BLOCK 243

% BEGIN SOURCE BLOCK 244 heading
\subsection{A deadline-count bound from eligibility}\label{sec:C.3}
% END SOURCE BLOCK 244

% BEGIN SOURCE BLOCK 245 paragraph
Let $d_r$ be the number of distinct deadline thresholds among eligible candidates at destination r. Starting at $E$, exact reconciliation reaches $b$ in at most $\min(n_r,2d_r-1)$ replacements per nonempty destination. Empty destinations need none. Nonnegative work makes $F$ order-reversing: $U\geq V$ implies $F(U)\leq F(V)$. Since $E\geq b$, odd iterates are lower bounds on $b$ and even iterates are upper bounds. We prove the stronger statement that any upper mask $b\leq U\leq E$ reaches $b$ within $2d-1$ replacements at one destination.
% END SOURCE BLOCK 245

% BEGIN SOURCE BLOCK 246 proof_open
\begin{proof}
% END SOURCE BLOCK 246

% BEGIN SOURCE BLOCK 247 paragraph
If $U=b$ there is nothing to prove. Otherwise let $p$ be its earliest extra admission: $U_p=1$, $b_p=0$. All earlier entries coincide with $b$, so after one replacement the entire prefix through $p$ is correct and remains correct. If $p$ is rejected by capacity, that fixed earlier prefix already fills capacity, and every later candidate is rejected; $F(U)=b$.
% END SOURCE BLOCK 247

% BEGIN SOURCE BLOCK 248 paragraph
Otherwise $p$ is gate-rejected and the fixed earlier workload is at least $D_p$. Every later candidate with $\text{deadline}\leq D_p$ must then be rejected, because subsequent work is nonnegative. These entries are zero after the first replacement and stay zero. After two replacements, $F^{2}(U)$ is again an upper bound on $b$, with that prefix and those later rejected classes fixed. Delete those fixed entries and absorb the fixed prefix into $W$ and Q. The remaining suffix has at most $d-1$ distinct deadlines, unchanged order and the same form of gate/capacity rule. Induction needs at most $2(d-1)-1$ further replacements, giving $2d-1$ overall. For $d=1$, the first gate rejection excludes every later candidate, so one replacement suffices. Capacity can only terminate the suffix earlier.
% END SOURCE BLOCK 248

% BEGIN SOURCE BLOCK 249 proof_close
\end{proof}
% END SOURCE BLOCK 249

% BEGIN SOURCE BLOCK 250 paragraph
Zero work is allowed: it need not advance workload, but every admission still advances count. Initial workload at/above a deadline, capacity already full, ineligible gaps and arbitrary ordering of the two thresholds do not defeat the proof. The same argument also gives one replacement when deadlines are nonincreasing in candidate order: the first extra gate rejection excludes every subsequent deadline class. The bound is sufficient; it need not be attained for a particular input.
% END SOURCE BLOCK 250

% BEGIN SOURCE BLOCK 251 paragraph
TrafficTwin's distinct production deadlines are 100 and 500 ms, so \textbf{three replacements equal the causal fixed-target mask in exact arithmetic} under these assumptions. Its service distributions are positive and do not weaken this conclusion; no independent choice of a convenient service/deadline is needed for the proof. Ingress with fixed individual targets also decomposes by destination. Causal per-task reselection does not have fixed targets, so this result does not equate A with C.
% END SOURCE BLOCK 251

% BEGIN SOURCE BLOCK 252 paragraph
The earlier five-task 1/1/41/41/81 example has three deadline thresholds and lies outside this two-threshold result. Its instability is therefore not evidence for an exact-arithmetic production-domain reconciliation defect. The new deduction explains why the earlier two-threshold grid was stable without using that grid as proof.
% END SOURCE BLOCK 252

% BEGIN SOURCE BLOCK 253 heading
\subsection{Why float32 requires a separate claim}\label{sec:C.4}
% END SOURCE BLOCK 253

% BEGIN SOURCE BLOCK 254 paragraph
The frozen helper computes an inclusive \texttt{jnp.cumsum(w, axis=0)} and subtracts w to form the exclusive prefix. Exact cancellation removes the current candidate; rounded addition/subtraction need not do so identically for both values of its own mask. It also adds base workload after prefix accumulation. The causal helper starts at base workload and adds admitted services in scan order; floating-point associativity differs. Thus the exact dependency and upper/lower arguments are not a proof of float32 equivalence or convergence.
% END SOURCE BLOCK 254

% BEGIN SOURCE BLOCK 255 paragraph
Count ranks use float32 cumulative sums too. Integer counts within the small suite and the studied slot widths are exactly representable, but that does not certify arbitrarily large populations. Strict gate comparisons require exact mask comparison, not an allclose acceptance criterion. The diagnostic protocol declares numerical reporting tolerances before execution, retains all discrete disagreements and separates task-parameter compatibility, whole-model reachability and observed frequency. No saved September candidate masks establish how often any constructed threshold case occurs.
% END SOURCE BLOCK 255

% BEGIN SOURCE BLOCK 256 heading
\subsection{Consequence for queue clearing in the stated exact model}\label{sec:C.5}
% END SOURCE BLOCK 256

% BEGIN SOURCE BLOCK 257 paragraph
With service multiplier 1, the largest possible RSU task service is less than 38.889 ms (type2, nominal 35.354 ms times noise below 1.1). At a destination with admissions, the last admitted task sees less than its deadline, at most 500 ms, before adding its own service. The final accumulated workload is therefore less than 538.889 ms, provided initial workload does not already exceed this bound. If nothing is admitted, workload is unchanged. Starting empty, induction across substeps preserves the bound for either fixed-target causal admission or causal reselection. The exact two-deadline result transfers it to three-pass fixed-target reconciliation.
% END SOURCE BLOCK 257

% BEGIN SOURCE BLOCK 258 paragraph
One 1,000 ms drain consequently clears every queue from this initialisation. Thus the empty-boundary assumption in the earlier destination-prefix explanation follows from the stated exact gate/service/timing model, not traffic sparsity alone. The observed zero endpoints and reconstructed maxima near 538 ms are consistent with that consequence. This is not a float32 error bound, a validation of event-level service or a theorem for other deadlines/service multipliers/drain intervals. Together with $R\geq K$, positive productive admissions and lowest-index ties, it supports recurrence of the low-index destination prefix; it does not identify every deadline loss.
% END SOURCE BLOCK 258

% BEGIN SOURCE BLOCK 259 heading
\subsection{Retained diagnostic evidence}\label{sec:C.6}
% END SOURCE BLOCK 259

% BEGIN SOURCE BLOCK 260 paragraph
The initial 671 inputs comprise 128 general mathematical fixtures, 324 coupled fixtures, 210 deliberate threshold fixtures and nine completion boundaries. Bounded adverse reductions/transfers and numerical follow-ups bring the total to 719, within the predeclared ceiling of 800. Protocol amendments precede their specific follow-ups; no outcome-driven expansion beyond that budget occurred. Full negative and unchanged results remain in the machine-readable record.
% END SOURCE BLOCK 260

% BEGIN SOURCE BLOCK 261 paragraph
The numerical alternating case uses seven type3 services of 2.5252525806427 ms followed by one type1 service of 21.11111068725586 ms, all with 100 ms deadlines. Initial $W=82.32322692871094$ ms and $Q=4$, $C=6{,}220$. Four earlier type1 services of $W/4$ construct this queue in a separate local admission substep; their noise multiplier lies within 0.9--1.1. This establishes local kernel compatibility under explicit arrays, not exact RNG/actor/traffic reachability.
% END SOURCE BLOCK 261

% BEGIN SOURCE BLOCK 262 paragraph
Exact arithmetic admits all eight, with the last gate margin about 0.000005007 ms; the last task nevertheless misses completion. The instrumented float32 mirror alternates $11111111\to11111110\to11111111\to11111110\to11111111$. The unchanged production helper confirms the retained third mask, recomputed offsets and capacity-failure count. Its causal float32 scan also rejects the last candidate, though via its own rounded gate. Thus this fixture distinguishes exact from implemented arithmetic and exposes an unstable label, not a demonstrated A\ensuremath{\to}B success improvement. The mirror's intermediate masks are constructed diagnostics, not recorded September fields.
% END SOURCE BLOCK 262

% BEGIN SOURCE BLOCK 263 paragraph
Four padding checks retain the alternating mask at $N=215/R=9$ and $N=2{,}488/R=10$, with the active block at either end. Their target is imposed: a busy target would not be the common argmin while others are empty. A final fixture prepends the four services, starts every queue empty and selects RSU0 at the correct $N=215/R=9$ dimensions. Its third mask is stable, although the last admission differs from exact arithmetic. This comparison prevents treating fixed-target instability as demonstrated on a full common-target trajectory.
% END SOURCE BLOCK 263

% BEGIN SOURCE BLOCK 264 paragraph
One separate inclusive-completion boundary also differs: saved-offset arithmetic rounds a just-late exact completion to an on-time float32 value. This is a zero-transfer calculation, not execution of the full latency scorer. Neither numerical example measures historical frequency or changes the reported fleet estimates.
% END SOURCE BLOCK 264

% BEGIN SOURCE BLOCK 265 paragraph
The two-RSU adverse case is minimal under the tested shared-position deletions and one-substep reduction. It retains production task/deadline/capacity values but changes infrastructure dimensions; its twelve transfers to the studied RSU/substep parameters retain no adverse result. No fixture frequency is interpreted as a fleet probability; that numerical investigation contains no confidence intervals, raw audits or full-system runs.
% END SOURCE BLOCK 265

% BEGIN SOURCE BLOCK 266 heading
\subsection{Retained adverse dispatch illustration}\label{sec:C.7}
% END SOURCE BLOCK 266

% BEGIN SOURCE BLOCK 267 paragraph
Target reselection can separately lose useful completions. Table~\ref{tab:C1} uses two empty RSUs, two substeps, capacity 6,220, zero transfer latency and repeated type2/type2/type1/type2 tasks with nominal production service demands. Both arms admit all eight tasks. B replays A's targets 0 then 1, whereas C chooses the current least workload after every reservation.
% END SOURCE BLOCK 267

% BEGIN SOURCE BLOCK 268 table
\begingroup
\small
\begin{xltabular}{\textwidth}{@{}YYYYYYY@{}}
\caption{Reduced-infrastructure illustration with production task/deadline/capacity values. Every listed task is admitted in both arms; completion is diagnostic service latency in milliseconds. ``Slot.task'' preserves candidate order. Values shown to three decimals; tests retain full precision.}\label{tab:C1} \\
\toprule
Slot.task & Deadline & Service & B: RSU / work before & B: completion & C: RSU / work before & C: completion \\
\midrule\endfirsthead
\toprule
Slot.task & Deadline & Service & B: RSU / work before & B: completion & C: RSU / work before & C: completion \\
\midrule\endhead
\bottomrule\endfoot
1.1 & 500 & 35.354 & 0 / 0.000 & 35.354 met & 0 / 0.000 & 35.354 met \\
1.2 & 500 & 35.354 & 0 / 35.354 & 70.707 met & 1 / 0.000 & 35.354 met \\
1.3 & 100 & 21.111 & 0 / 70.707 & 91.818 met & 0 / 35.354 & 56.465 met \\
1.4 & 500 & 35.354 & 0 / 91.818 & 127.172 met & 1 / 35.354 & 70.707 met \\
2.1 & 500 & 35.354 & 1 / 0.000 & 35.354 met & 0 / 56.465 & 91.818 met \\
2.2 & 500 & 35.354 & 1 / 35.354 & 70.707 met & 1 / 70.707 & 106.061 met \\
2.3 & 100 & 21.111 & 1 / 70.707 & 91.818 met & 0 / 91.818 & 112.929 \textbf{miss} \\
2.4 & 500 & 35.354 & 1 / 91.818 & 127.172 met & 1 / 106.061 & 141.414 met \\
\end{xltabular}
\endgroup
% END SOURCE BLOCK 268

% BEGIN SOURCE BLOCK 269 paragraph
B preserves an empty RSU for substep 2. C spreads earlier work across both, making the later 100 ms task finish at 112.929 ms: eight admissions produce seven successes. Current workload omits future deadlines, and the gate omits own service; neither directly maximises deadline-success count. Tested deletions remove this loss, without establishing global minimality.
% END SOURCE BLOCK 269

% BEGIN SOURCE BLOCK 270 paragraph
The loss did not persist in twelve tested transfers to nine/ten RSUs and five substeps, with four/six candidates. The example identifies an order/service/deadline interaction, not the cause of recorded losses at full dimensions. B replays A targets and is not an autonomous deployable comparator. Original full-evaluator prefixes remain unrun \hyperref[source-s10]{[S10]}, \hyperref[source-s14]{[S14]}.
% END SOURCE BLOCK 270

% BEGIN SOURCE BLOCK 271 heading
\appendixsection{app:D}{Type aggregation and engineering evidence}
% END SOURCE BLOCK 271

% BEGIN SOURCE BLOCK 272 paragraph
The retained retrospective type analysis read eight original primary-morning task arrays, authenticated their hashes, and aggregated categories by type; it was not repeated here. The prior 19-run admission validation and 23 task joins are reused, not repeated. All-task totals and paired net changes reconcile. The following table exposes every type/draw/arm; attainment uses the offered column. Other terminal failures are available in the linked CSV and unchanged within each pair.
% END SOURCE BLOCK 272

% BEGIN SOURCE BLOCK 273 table
\begingroup
\small
\begin{xltabular}{\textwidth}{@{}YYYYYYYY@{}}
\caption{Type outcomes. I = ingress, P = per-task; offered and category counts are tasks, attainment is percent. Gate rejection and admitted misses are disjoint.}\label{tab:D1} \\
\toprule
Seed / type & Arm & Offered & Success & Attainment & Admitted & Gate rejected & Admitted miss \\
\midrule\endfirsthead
\toprule
Seed / type & Arm & Offered & Success & Attainment & Admitted & Gate rejected & Admitted miss \\
\midrule\endhead
\bottomrule\endfoot
0/1 & I & 349,045 & 252,283 & 72.278 & 328,923 & 19,548 & 76,640 \\
0/1 & P & 349,045 & 272,106 & 77.957 & 347,605 & 866 & 75,499 \\
0/2 & I & 523,060 & 516,034 & 98.657 & 522,134 & 46 & 6,100 \\
0/2 & P & 523,060 & 516,106 & 98.671 & 522,180 & 0 & 6,074 \\
0/3 & I & 872,011 & 780,747 & 89.534 & 821,868 & 48,572 & 41,121 \\
0/3 & P & 872,011 & 828,664 & 95.029 & 868,284 & 2,156 & 39,620 \\
2/1 & I & 349,045 & 242,629 & 69.512 & 324,165 & 24,149 & 81,536 \\
2/1 & P & 349,045 & 261,794 & 75.003 & 347,317 & 997 & 85,523 \\
2/2 & I & 523,060 & 515,409 & 98.537 & 521,901 & 40 & 6,492 \\
2/2 & P & 523,060 & 515,471 & 98.549 & 521,941 & 0 & 6,470 \\
2/3 & I & 872,011 & 768,687 & 88.151 & 809,915 & 60,362 & 41,228 \\
2/3 & P & 872,011 & 828,363 & 94.995 & 867,771 & 2,506 & 39,408 \\
3/1 & I & 349,045 & 253,921 & 72.747 & 330,364 & 18,383 & 76,443 \\
3/1 & P & 349,045 & 274,260 & 78.574 & 348,537 & 210 & 74,277 \\
3/2 & I & 523,060 & 517,080 & 98.857 & 522,533 & 1 & 5,453 \\
3/2 & P & 523,060 & 517,086 & 98.858 & 522,534 & 0 & 5,448 \\
3/3 & I & 872,011 & 782,701 & 89.758 & 824,760 & 46,357 & 42,059 \\
3/3 & P & 872,011 & 830,305 & 95.217 & 870,539 & 578 & 40,234 \\
4/1 & I & 349,045 & 260,449 & 74.618 & 333,800 & 14,596 & 73,351 \\
4/1 & P & 349,045 & 280,703 & 80.420 & 348,250 & 146 & 67,547 \\
4/2 & I & 523,060 & 516,064 & 98.662 & 521,928 & 1 & 5,864 \\
4/2 & P & 523,060 & 516,069 & 98.663 & 521,929 & 0 & 5,860 \\
4/3 & I & 872,011 & 795,155 & 91.186 & 834,388 & 35,739 & 39,233 \\
4/3 & P & 872,011 & 832,211 & 95.436 & 869,818 & 309 & 37,607 \\
\end{xltabular}
\endgroup
% END SOURCE BLOCK 273

% BEGIN SOURCE BLOCK 274 paragraph
The \href{\detokenize{../empirical_extension_2026-09-08/analysis/results/BENCHMARK.json}}{benchmark record} retains lowering/compilation times, median, quartiles, p95, minima/maxima, compiler buffer estimates and exact source/runtime identities for all 16 configurations. \href{\detokenize{../empirical_extension_2026-09-08/analysis/results/BENCHMARK_TIMES.csv}}{Raw call durations} permit descriptive timing checks. \href{\detokenize{../empirical_extension_2026-09-08/confirmation/PRECISION.json}}{Precision sensitivity} records pre-outcome assumptions about block variance; it is not observed power for the new results.
% END SOURCE BLOCK 274

% BEGIN SOURCE BLOCK 275 heading
\appendixsection{app:E}{Joint-randomness study controls and execution}
% END SOURCE BLOCK 275

% BEGIN SOURCE BLOCK 276 paragraph
The execution-authorised seal preserves the disabled preparation package and binds its own source version, current actor/trace identities, runtime, output root, lock and qualification receipts. Fleet/evaluator pairs are (100,200) through (107,207). The accessible manifest check found no prior outcomes for these pairs; this is a scoped absence check, not proof that no record exists elsewhere. Block 0 starts ingress/common-target/per-task/round-robin; each subsequent block rotates that order by one position.
% END SOURCE BLOCK 276

% BEGIN SOURCE BLOCK 277 paragraph
Qualification used the previously inspected fleet/evaluator pair (1,0), never confirmation outcomes. Three frozen/instrumented pairs ran for 300 steps, followed by round-robin at 300 and a fresh-process 150-step restart: eight attempts total. All 83 shared scientific fields matched exactly in each unchanged-arm pair. The 300-step cyclic run advanced its pointer 8,817 times, including 28 subsequently rejected attempts; its final pointer was 6. The 150-step run matched the longer run's prefix. No unavailable-radio, RSU-capacity or local-capacity failures occurred in these short records, so their end-to-end branches were not empirically exercised. Preserved kernel tests cover those branches separately. Qualification is neither an extra replication nor proof of full-horizon equivalence.
% END SOURCE BLOCK 277

% BEGIN SOURCE BLOCK 278 paragraph
Added records expose final admission independently of success, actual local/peer service, queue counts, post-service workloads, SoC/energy and cyclic pointer state. Negative checks refused a corrupted offered count, mismatched shared-input binding, absent block receipt and changed output hash. A separate prelaunch critique identified a possible alternate-root attempt-budget bypass, incomplete type-summary checking and misleading failed-process timing. The final version binds one raw root/global lock, validates type numerators/denominators and records failure phase and elapsed process time. Existing short records were rechecked; no ninth qualification evaluator ran.
% END SOURCE BLOCK 278

% BEGIN SOURCE BLOCK 279 paragraph
Each full cell requires exact discrete accounting, proposal/admission/execution/forwarding consistency, penalty latency exactly ten times the deadline, inclusive deadline success and independent queue-work reconstruction. Retained tolerances are 0.01 ms absolute per queue-second for service conservation, 0.001 ms plus relative 0.000001 for drain, and 0.000001 task for reconstructed JSON score/type numerators. Within each block, fleet arrays, key streams, offers, observation descriptors, operational types/sizes and pre-admission service work must agree exactly. Observations, logits, actions, SoC and selected radio quantities are separately reported; equality is not forced.
% END SOURCE BLOCK 279

% BEGIN SOURCE BLOCK 280 paragraph
Execution is serial with at most 32 full attempts including failures, no retries or seed substitution, a three-hour per-cell timeout, 50 GiB initial free-space requirement and 20 GiB before each attempt. Incomplete attempts are retained. Completed-output hashes, commands, source identities and all eight block receipts gate the sealed primary analysis. The first included block supplies a remaining-runtime estimate; effects cannot change continuation. Process time includes startup, compilation, evaluation and compressed output writing; validation is timed separately. These times are not the scheduler-only benchmark in Table~\ref{tab:10}.
% END SOURCE BLOCK 280

% BEGIN SOURCE BLOCK 281 paragraph
The completed full matrix has 32 valid cells, eight valid block receipts and no failed attempts. The following tables retain every cell and block. I = ingress, $D$ = common-target, P = per-task and $R$ = round-robin; fleet/evaluator seeds are 100+$b$ and 200+$b$ for block b. Other terminal categories are retained in CELL\_RESULTS.csv and each validation receipt.
% END SOURCE BLOCK 281

% BEGIN SOURCE BLOCK 282 table
\begingroup
\small
\begin{xltabular}{\textwidth}{@{}YYYYYYY@{}}
\caption{Full new-study counts. Offered populations match within blocks; attainment is offered-task percent. Gate rejection and admitted misses are disjoint.}\label{tab:E1} \\
\toprule
Block / arm & Offered & Admitted & Successes & Attainment & Gate rejected & Admitted misses \\
\midrule\endfirsthead
\toprule
Block / arm & Offered & Admitted & Successes & Attainment & Gate rejected & Admitted misses \\
\midrule\endhead
\bottomrule\endfoot
0/I & 1,744,761 & 1,655,779 & 1,525,913 & 87.457 & 86,625 & 129,866 \\
0/$D$ & 1,744,761 & 1,582,636 & 1,459,384 & 83.644 & 159,768 & 123,252 \\
0/P & 1,744,761 & 1,737,398 & 1,605,152 & 91.998 & 5,006 & 132,246 \\
0/$R$ & 1,744,761 & 1,726,560 & 1,592,830 & 91.292 & 15,844 & 133,730 \\
1/I & 1,740,970 & 1,662,805 & 1,534,040 & 88.114 & 76,351 & 128,765 \\
1/$D$ & 1,740,970 & 1,593,599 & 1,470,819 & 84.483 & 145,557 & 122,780 \\
1/P & 1,740,970 & 1,736,301 & 1,607,919 & 92.358 & 2,855 & 128,382 \\
1/$R$ & 1,740,970 & 1,727,104 & 1,596,377 & 91.695 & 12,052 & 130,727 \\
2/I & 1,741,855 & 1,639,136 & 1,510,719 & 86.730 & 99,922 & 128,417 \\
2/$D$ & 1,741,855 & 1,556,771 & 1,436,166 & 82.450 & 182,287 & 120,605 \\
2/P & 1,741,855 & 1,732,694 & 1,598,164 & 91.751 & 6,364 & 134,530 \\
2/$R$ & 1,741,855 & 1,718,789 & 1,583,772 & 90.924 & 20,269 & 135,017 \\
3/I & 1,744,227 & 1,677,045 & 1,557,192 & 89.277 & 65,368 & 119,853 \\
3/$D$ & 1,744,227 & 1,616,965 & 1,501,691 & 86.095 & 125,448 & 115,274 \\
3/P & 1,744,227 & 1,740,631 & 1,624,157 & 93.116 & 1,782 & 116,474 \\
3/$R$ & 1,744,227 & 1,734,468 & 1,614,712 & 92.575 & 7,945 & 119,756 \\
4/I & 1,743,484 & 1,680,263 & 1,560,221 & 89.489 & 61,226 & 120,042 \\
4/$D$ & 1,743,484 & 1,623,533 & 1,507,624 & 86.472 & 117,956 & 115,909 \\
4/P & 1,743,484 & 1,738,717 & 1,622,888 & 93.083 & 2,772 & 115,829 \\
4/$R$ & 1,743,484 & 1,732,685 & 1,613,895 & 92.567 & 8,804 & 118,790 \\
5/I & 1,745,078 & 1,685,944 & 1,564,006 & 89.624 & 56,788 & 121,938 \\
5/$D$ & 1,745,078 & 1,629,177 & 1,511,177 & 86.597 & 113,555 & 118,000 \\
5/P & 1,745,078 & 1,740,768 & 1,623,996 & 93.062 & 1,964 & 116,772 \\
5/$R$ & 1,745,078 & 1,735,213 & 1,614,934 & 92.542 & 7,519 & 120,279 \\
6/I & 1,742,095 & 1,669,339 & 1,541,879 & 88.507 & 69,142 & 127,460 \\
6/$D$ & 1,742,095 & 1,603,847 & 1,481,655 & 85.050 & 134,634 & 122,192 \\
6/P & 1,742,095 & 1,736,400 & 1,611,494 & 92.503 & 2,081 & 124,906 \\
6/$R$ & 1,742,095 & 1,729,087 & 1,601,138 & 91.909 & 9,394 & 127,949 \\
7/I & 1,744,742 & 1,657,846 & 1,531,566 & 87.782 & 84,523 & 126,280 \\
7/$D$ & 1,744,742 & 1,588,286 & 1,468,334 & 84.158 & 154,083 & 119,952 \\
7/P & 1,744,742 & 1,736,713 & 1,608,728 & 92.204 & 5,656 & 127,985 \\
7/$R$ & 1,744,742 & 1,726,431 & 1,596,904 & 91.527 & 15,938 & 129,527 \\
\end{xltabular}
\endgroup
% END SOURCE BLOCK 282

% BEGIN SOURCE BLOCK 283 table
\begingroup
\small
\begin{xltabular}{\textwidth}{@{}YYYY@{}}
\caption{All eight paired effects, percentage points. These blocks, not tasks or RSUs, are the new replications.}\label{tab:E2} \\
\toprule
Block & Per-task \ensuremath{-} ingress & Common-target \ensuremath{-} ingress & Per-task \ensuremath{-} round-robin \\
\midrule\endfirsthead
\toprule
Block & Per-task \ensuremath{-} ingress & Common-target \ensuremath{-} ingress & Per-task \ensuremath{-} round-robin \\
\midrule\endhead
\bottomrule\endfoot
0 & +4.542 & \ensuremath{-}3.813 & +0.706 \\
1 & +4.244 & \ensuremath{-}3.631 & +0.663 \\
2 & +5.020 & \ensuremath{-}4.280 & +0.826 \\
3 & +3.839 & \ensuremath{-}3.182 & +0.542 \\
4 & +3.594 & \ensuremath{-}3.017 & +0.516 \\
5 & +3.438 & \ensuremath{-}3.027 & +0.519 \\
6 & +3.996 & \ensuremath{-}3.457 & +0.594 \\
7 & +4.423 & \ensuremath{-}3.624 & +0.678 \\
\end{xltabular}
\endgroup
% END SOURCE BLOCK 283

% BEGIN SOURCE BLOCK 284 paragraph
Actions matched exactly in all 24 non-reference arm/block comparisons. Small observation and logit differences occurred in block 1; they did not change actions. The estimand remains the total scheduling intervention under frozen weights. Observed action matching extends the original finding to these streams without forcing replay or guaranteeing future equality. In block 1 (fleet 101/evaluator 201), each alternative differed from ingress by at most 0.0000000149012 in observations and SoC, and 0.000000357628 in logits. These are dimensionless recorded differences; no action was altered or forced. The \href{\detokenize{../joint_confirmation_2026-09-08/evidence/BLOCK_CONTROLS.json}}{block receipts} retain every comparison and exact maxima. The \href{\detokenize{../joint_confirmation_2026-09-08/evidence/CELL_RESULTS.json}}{cell records} and \href{\detokenize{../joint_confirmation_2026-09-08/evidence/PAIRED_EFFECTS.csv}}{paired effects} preserve unrounded values. No task-level significance tests or new type subgroup analysis were added.
% END SOURCE BLOCK 284

% BEGIN SOURCE BLOCK 285 paragraph
All 32 records passed final-admission/score/category/penalty, assignment and queue checks under the sealed contract. Maximum reconstructed per-queue-second service discrepancy was 0.000121593 ms for RSUs and 0.000244141 ms for vehicle queues. This finite result does not certify every possible path or repair the historical off scorer.
% END SOURCE BLOCK 285

% BEGIN SOURCE BLOCK 286 paragraph
The campaign elapsed 6440.25 seconds (107.34 minutes) from first attempt to completion. Full evaluator processes took 6375.54 seconds in total; cell validation took 49.59 seconds, block validation 10.83 seconds, and sealed analysis with completion/hash checks 3.44 seconds. Qualification processes separately took 68.09 seconds. \href{\detokenize{../joint_confirmation_2026-09-08/evidence/TIMING.json}}{Timing records} and per-cell logs retain the exact boundaries. The first-block estimate used runtime only.
% END SOURCE BLOCK 286

% BEGIN SOURCE BLOCK 287 paragraph
The \href{\detokenize{../joint_confirmation_2026-09-08/evidence/RAW_INVENTORY.json}}{raw inventory} binds the locally retained task/step arrays and supporting files. The \href{\detokenize{../joint_confirmation_2026-09-08/document/verify_results.py}}{portable verifier} regenerates central tables from compact evidence and accepts an optional explicit raw root for integrity checking. Compact arithmetic is not repeated task-level validation; hashes are not a backup. Approved off-machine storage and examiner access remain unresolved.
% END SOURCE BLOCK 287

\end{document}
